% !TeX program = latexmk
% !TeX encoding = UTF-8
% 请使用 paper-Chinese/build.sh 完整编译；不要仅运行一次 xelatex。
\documentclass{ieeeaccess}
\usepackage{cite}
\usepackage{amsmath,amssymb,amsfonts}
\usepackage{algorithmic}
\usepackage{graphicx}
\usepackage{textcomp}
\usepackage{booktabs}
\usepackage{multirow}
\usepackage{xcolor}
\usepackage{url}

% 中文支持（中文草稿阶段使用 XeLaTeX 编译）
\usepackage{xeCJK}
\setCJKmainfont{STSong}
\setCJKsansfont{STHeiti}
\setCJKmonofont{STFangsong}

\usepackage{bm}
\makeatletter
\AtBeginDocument{\DeclareMathVersion{bold}
\SetSymbolFont{operators}{bold}{T1}{times}{b}{n}
\SetSymbolFont{NewLetters}{bold}{T1}{times}{b}{it}
\SetMathAlphabet{\mathrm}{bold}{T1}{times}{b}{n}
\SetMathAlphabet{\mathit}{bold}{T1}{times}{b}{it}
\SetMathAlphabet{\mathbf}{bold}{T1}{times}{b}{n}
\SetMathAlphabet{\mathtt}{bold}{OT1}{pcr}{b}{n}
\SetSymbolFont{symbols}{bold}{OMS}{cmsy}{b}{n}
\renewcommand\boldmath{\@nomath\boldmath\mathversion{bold}}}
\makeatother

\def\BibTeX{{\rm B\kern-.05em{\sc i\kern-.025em b}\kern-.08em
    T\kern-.1667em\lower.7ex\hbox{E}\kern-.125emX}}

% 投稿占位符有意不标红，但仍需在视觉上明确显示。
% 生成终稿前请全文检索“占位”。
\newcommand{\TODO}[1]{\PackageWarning{manuscript}{Unresolved submission placeholder}\textbf{[占位：#1]}}

%% ============================================================
\begin{document}
\history{出版日期 xxxx 00, 0000，当前版本日期 xxxx 00, 0000。}
\doi{10.1109/ACCESS.2026.DOI}

\title{标称 3.91 ms UVC 设置下的时序像素掩蔽：一项传统 OCR 档位级测量研究}

\author{\uppercase{First A. Author}\authorrefmark{1}
\TODO{作者信息待填}}

\address[1]{\TODO{单位及地址待填}}

\tfootnote{\TODO{基金致谢待填。}
本文包含一项计划开展的最小风险人类被试研究。目前尚未
采集任何被试数据。只有在主管机构或院系出具
书面批准、豁免或免审认定后，方可开始招募；计划采用的保护措施
见第~\ref{sec:user_study} 节。}

\markboth
{作者 \headeretal：标称 UVC 设置下的时序像素掩蔽}
{作者 \headeretal：标称 UVC 设置下的时序像素掩蔽}

\corresp{通讯作者：\TODO{通讯作者邮箱}。}

%% ============================================================
%% 摘要
%% ============================================================
\begin{abstract}
屏幕内容可被静默拍摄，继而由光学字符识别（OCR）或视觉语言模型（VLM）处理。本研究测量在一个标称 USB 视频类（UVC）曝光控制设置下，传统 OCR 恢复率随档位发生的变化。该处理管线将每帧划分为高速交替的互补子帧，并加入随档位变化的叠加层。我们在 9 组距离--角度配置下共采集 10,575 张真机拍摄图像。主要的共同设置分析对三个档位分别使用同一组 288 个经重复采集平均的短曝光单元（12 个内容项 $\times$ 8 个拍摄几何条件 $\times$ 3 次重复）。按单次采集计算的引擎最优平均字符恢复率在未保护档、可读性优先档和高抑制档下分别为 94.5\%、17.8\% 和 5.6\%。在 432 个显式标注的字段评估机会中，可读性优先档仍有 17.7\% 的字段微平均完整恢复率，其匹配 P99 字符恢复率为 95.0\%。固定五变换、三引擎的网格最优选择（oracle）将可读性优先档和高抑制档的恢复率分别提高到 40.2\% 和 13.7\%。150 帧时域均值攻击和正轴 VLM 探针恢复出了多得多的文本。本实验没有确定有效曝光区间，也未将时序稀疏性与叠加层及占空比效应分离，因而不支持跨设备保护主张。本文的贡献是在单一 UVC 链路下获得的一项有明确边界的实证测量，以及足以排除部署主张或通用屏幕防拍摄主张的失效证据。
\end{abstract}

\begin{keywords}
对抗扰动，相机拍摄，显示隐私，光学字符识别，视觉暂留，屏幕--相机拍摄，时序掩蔽，视觉窃听
\end{keywords}

\titlepgskip=-21pt
\maketitle

%% ============================================================
%% I. 引言
%% ============================================================
\section{引言}

\PARstart{屏}{幕}遍布办公室、银行、医院和会议室；拍摄设备也同样随处可见，包括智能手机、智能眼镜~\cite{b_wang2026}（例如 Ray-Ban Meta）以及微型相机。仅凭一张屏幕照片，攻击者便可通过光学字符识别（OCR）提取文字，或通过目标检测识别界面元素。这种\emph{视觉窃听}十分隐蔽：它既不留下网络痕迹，也无需与目标设备发生物理接触。受控实验与实地研究均证实，此类攻击较为常见、成功率高，而且往往不会被察觉~\cite{b_ponemon2016,b_eiband2017}。尽管上述设备构成了本文的研究动机，但这里给出的物理实验证据仅覆盖一款 USB 视频类（UVC）网络摄像头；智能手机、智能眼镜及其他图像信号处理器（ISP）管线上的可复现性仍有待未来研究。

\subsection{现有方法及其局限}

现有屏幕隐私解决方案各自仅适用于有限场景。空间、软件、水印和时序方法分别针对不同的观察几何条件或攻击阶段。已有文献尚未报告与本文相同的场景组合，即正轴短曝光拍摄、静态敏感文本以及拍摄后的机器识别；详细比较见 \S\ref{sec:related}。因此，有必要针对机器从真机拍摄快照中恢复信息的能力开展专门测量。

难点在于，对相机端信息的削弱往往也会损害合法观看者的可读性。本文考察显示的时间结构能否在短曝光拍摄下扩大二者之间的差距。

\subsection{核心思想}

本文方法利用人类时间整合与相机曝光窗口之间的差异。每一帧被分解为 $n$ 个随机互补的子帧，并以高速交替显示。当曝光时间短于或接近一个子帧的持续时间时，相机只会记录不完整的像素子集，而人类观察者则会跨子帧整合内容（图~\ref{fig:concept}）。由于相机曝光并非瞬时采样，覆盖完整周期的长曝光或视频重构能够重新累加这些互补子帧；因此，本文不对任意曝光条件作不可恢复性保证。

本文将工作定位为受控测量研究，而非部署主张。我们考察在一个标称 3.91\,ms UVC 控制设置下，档位级显示处理如何改变传统 OCR 的恢复率。其操作场景是一条仪器化管线：每个目标由一张归档图像表示，随后再由传统 OCR 处理。该场景并非普通智能手机拍摄、智能眼镜拍摄，也不涉及能够针对性优化拍摄与识别过程的操作者。在 8 个共同设置的拍摄几何条件下，匹配平均恢复率在未保护档、可读性优先档和高抑制档下分别为 94.5\%、17.8\% 和 5.6\%。高抑制档未达到其描述性的 $<5\%$ 目标。可以成立的主张仅限于单一 UVC 条件下的这种档位级关联，而不是某个有效曝光区间、通用防拍摄设置，或时序稀疏性单独导致恢复率下降的证明。

边界证据本身也是测量贡献的一部分，而非附带说明。实验未记录环境照度，因而无法将日志中的 UVC 控制设置与典型智能手机自动曝光条件作强弱排序。所评估的时域均值攻击对 150 帧短曝光图像取平均，按 60\,fps 计算名义时长为 2.5\,s；本文未评估更短的视频持续时间。视觉语言模型（VLM）探针能恢复相当一部分大字号短片段，这表明对传统 OCR 的抑制无法迁移至现代端到端识别器（\S\ref{sec:vlm}）。

\begin{figure}[!t]
\centering
\includegraphics[width=\columnwidth,viewport=0 0 252 168,clip]{figures/visio/figure1_concept_threat/final/figure1_concept_threat.pdf}
\caption{时序像素掩蔽的威胁场景与运行边界。(a) 附近的相机无需访问主机系统，即可拍摄屏幕上的敏感内容。(b) 人类视觉会整合互补子帧，而相机的一次短曝光可能只与显示周期的一部分重叠，留给 OCR 的只是一个不完整片段。图中的相机路径仅为概念示意：它不假设零时长采样，也不假设固定的人类整合时间；长曝光或多帧重构能够恢复多得多的信息。}
\label{fig:concept}
\end{figure}

\subsection{贡献}

本文的主要贡献如下：
\begin{enumerate}
\item \textbf{受控的档位级物理测量}：在既有时序掩蔽思想的基础上，我们实现了像素级互补子帧、随档位变化的条纹/字形叠加层以及图形处理器（GPU）渲染原型。归档数据包含使用一款 eMeet SmartCam S600、在 3 个距离和 3 个角度下采集的 10,575 张图像。每个档位包含 288 个共同设置匹配单元（12 个内容项 $\times$ 8 个拍摄几何条件 $\times$ 3 次重复），其中未保护档、可读性优先档和高抑制档的平均恢复率分别为 94.5\%、17.8\% 和 5.6\%。这是在固定条件下，对经过几何校正的裁剪图像使用三个开源 OCR 引擎处理后得到的观察结果。

\item \textbf{考虑聚类的描述性证据与曝光敏感性}：我们报告采集层面的分布及按内容项聚类的配对对比。在归档数据的全部 9 个拍摄几何条件下，未保护短曝光相较于可读性优先档高出 78.0 个百分点（内容项聚类重采样区间 [75.5, 80.3]），相较于高抑制档高出 89.1 个百分点（[85.0, 92.4]）。将分析限制于 8 个共同设置的拍摄几何条件后，相应差值分别为 76.7 个百分点（[74.3, 79.2]）和 88.9 个百分点（[85.5, 91.9]）。这些区间描述结果对 12 个归档内容聚类的敏感性，并将所测试的拍摄几何条件视为固定；它们并不量化跨设备或跨场景的不确定性。

\item \textbf{失效边界刻画}：逐样本最强数字攻击的恢复率达到 95.2\%。对 150 帧相机图像取均值后，可读性优先档和高抑制档的恢复率分别达到 71.1\% 和 47.9\%；排除欠曝光的拍摄几何条件后，又分别升至 79.7\% 和 53.9\%。VLM 探针可恢复相当一部分大字号内容。这些负面结果界定了上述狭义传统 OCR 观察的终点：可读性优先档是唯一与计划开展的可用性协议相关联的档位，而高抑制档仍只是一个面向安全性的探索性压力测试档位，尚无用户可接受性证据。
\end{enumerate}

本文余下内容安排如下。\S\ref{sec:related} 回顾相关工作，\S\ref{sec:threat} 定义威胁模型，\S\ref{sec:method} 介绍系统设计。\S\ref{sec:evaluation} 报告主要 OCR 证据、VLM 与积分边界，以及延期执行的可用性协议。\S\S\ref{sec:discussion} 和~\ref{sec:conclusion} 综合证据层级并总结全文。检测/跟踪诊断被置于补充材料，因为它们并未扩展主要 OCR 主张。


%% ============================================================
%% II. 相关工作
%% ============================================================
\section{相关工作}
\label{sec:related}

\subsection{屏幕视觉窃听威胁}

已有研究记录了多种攻击，包括远距离屏幕读取攻击（望远镜、长焦镜头）~\cite{b_backes2009}、对 HDMI 电磁泄漏的深度学习重构~\cite{b_fernandez2024}、智能眼镜隐蔽拍摄~\cite{b_wang2026}，以及通过反射和折射实施的间接窃听~\cite{b_raguram2011}。Ponemon Institute 的一项受控研究表明，视觉入侵尝试的总体成功率为 91\%，其中 49\% 的情形在 15~分钟内出现首次成功入侵~\cite{b_ponemon2016}。Eiband 等~\cite{b_eiband2017} 调查了 174 名参与者，以研究肩窥的普遍程度及用户采用的应对策略。Tang 和 Shin~\cite{b_tang2023} 提出了 Eye-Shield：该方法利用前置摄像头检测旁观者，并实时模糊敏感区域，但它依赖用户侧硬件，且只能提供局部掩蔽。本文关注无需攻击者侧配合的显示侧缓解措施，但不声称能够在所有曝光和重构条件下阻止信息获取。

\subsection{显示隐私保护}

物理隐私滤镜能够限制可视角度，但无法阻挡正轴相机。Chen 等\ 提出 HideScreen，这是一种显示侧肩窥防御方案，通过抑制低频视觉线索，使屏幕内容在超出设计观看距离后与背景融为一体~\cite{b_hidescreen2019}；它针对空间域中的人类肩窥，而非短曝光相机拍摄或机器识别。Zhu 等~\cite{b_zhu2017} 提出了 LiShield，通过调制智能 LED 照明来干扰滚动快门相机，但该方案作用于环境照明，而非显示内容。视觉密码术将秘密图像拆分为多个份额，并在叠加时将其组合~\cite{b_naor1994}，由此启发了时序分解思想，但其原始形式面向印刷介质。数字水印~\cite{b_zhong2023} 可嵌入追踪标识，其中包括抗屏摄变体~\cite{b_fang2019}，但它提供的是事后归因，而非事前阻断。基于软件的数字版权管理（DRM）能够阻止数字截图，却无法防止物理拍摄。

屏幕--相机可见光通信（VLC）~\cite{b_nguyen2016,b_tran2021,b_ghiasi2024} 同样利用人类时间整合，将人眼不可见但可由相机解码的信息嵌入显示内容。本文方法与 VLC 互为镜像：VLC 旨在让相机\emph{读取}隐藏信息，而本文则旨在降低受限短曝光拍摄下显示内容对机器的可读性。

Yerazunis 和 Carbone 的早期工作~\cite{b_yerazunis2001} 探讨了用于显示隐私的时序图像掩蔽。Wu 和 Zhai~\cite{b_wu2013} 对时序心理视觉调制进行了系统化研究，Gao 等~\cite{b_gao2014} 则将其用于信息安全显示。Zhang 等~\cite{b_kaleido2015} 提出了 Kaleido，通过多帧分解对视频进行重新编码。Rainbow~\cite{b_rainbow2018} 则调制环境照明，以破坏滚动快门相机对实体物体的拍摄。这些方案使用的成像指标和主观质量指标，无法与本文的恢复率指标互换。既有研究没有报告本文所考察的特定组合：静态文本、短 UVC 曝光以及拍摄后的识别。Kaleido 使用色度互补帧与亮度变化，而本文的实现采用随机离散像素划分。探索性的高抑制档还放弃了严格完备性（\S\ref{sec:hardened}）。Lim~\cite{b_lim2007} 利用视觉暂留在浏览器中生成防屏幕截取键盘。表~\ref{tab:kaleido_compare} 将本研究置于有代表性的空间、照明和时序方案之中加以比较。

\begin{table*}[!t]
\caption{与代表性显示及实体场景隐私方案的适用范围比较。表中各方案针对不同威胁与评估目标，不构成参数匹配的性能主张。}
\label{tab:kaleido_compare}
\centering
\footnotesize
\setlength{\tabcolsep}{3.2pt}
\begin{tabular}{p{0.15\textwidth}p{0.15\textwidth}p{0.20\textwidth}p{0.25\textwidth}p{0.17\textwidth}}
\toprule
方案 & 保护介质 & 目标威胁 & 机制 & 已报告评估 \\
\midrule
HideScreen~\cite{b_hidescreen2019} & 移动显示屏 & 超出设计距离的人类肩窥 & 空间域低频线索抑制 & 人类观看距离/可读性 \\
LiShield~\cite{b_zhu2017} & 实体场景 & 滚动快门相机拍摄 & 调制环境 LED 照明 & 拍摄图像干扰 \\
Kaleido~\cite{b_kaleido2015} & 显示视频 & 连续录像盗录 & 感知完备的色度/亮度重编码 & 峰值信噪比（PSNR）、结构相似性指数（SSIM）和主观质量 \\
Rainbow~\cite{b_rainbow2018} & 实体物体 & 移动相机盗录 & 环境照明调制 & 滚动快门图像退化 \\
Lim~\cite{b_lim2007} & 浏览器键盘 & 单次数字屏幕截取 & 视觉暂留局部图案 & 键盘截取演示 \\
本文 & 静态屏幕文本 & 短曝光物理拍摄后的机器识别 & 随机像素划分；高抑制档破坏完备性 & OCR/VLM 恢复率与积分边界 \\
\bottomrule
\end{tabular}
\end{table*}

本文的贡献并非孤立的时序划分，而是在单一 UVC 条件下开展的系统级测量，并同时给出显式字段、积分及 VLM 失效证据。检测、跟踪和数字可用性代理仍属于补充诊断。

\subsection{对抗扰动与物理世界鲁棒性}

Wei 等~\cite{b_wei2024} 综述了数字对抗扰动~\cite{b_goodfellow2015,b_madry2018,b_carlini2017} 与物理世界对抗补丁~\cite{b_eykholt2018,b_thys2019,b_xu2020}。Athalye 等~\cite{b_athalye2018} 通过变换期望（Expectation over Transformation, EoT）首次实现了跨视角鲁棒的 3D 物理对抗物体，但该方法仍需对白盒目标模型具有访问权限。物理部署会引入数字评估中不存在的变换，因此必须在不同拍摄条件下测试鲁棒性，而不能从像素空间性能直接推断~\cite{b_eykholt2018}。物理对抗补丁具有跨视角鲁棒性，但需要针对特定分类器预先生成扰动图案~\cite{b_thys2019,b_xu2020}，对未见模型的迁移能力有限~\cite{b_gu2024}。屏幕--相机链路还会引入摩尔纹（moir\'{e}）伪影；Niu 等~\cite{b_niu2021} 已证明，这类伪影本身也可充当对抗扰动。基础机制在显示层对内容进行时序划分，而非优化模型特定的空间扰动。增强档位还会加入噪声及条纹/字形叠加层，其物理迁移效果由实验评估，而非预先假定。物理摄影本质上是一个采样过程，而防护机制则嵌入显示本身。


%% ============================================================
%% III. 威胁模型
%% ============================================================
\section{威胁模型}
\label{sec:threat}

\subsection{受保护资产与评估内容}

研究动机所涉及的资产包括密码、金融数据、个人身份信息（PII）、源代码和机密文件。实验并未表明所有此类资产都能得到保护。语料覆盖凭据、数字串、URL、代码、段落以及一个文档页面。我们在评分前对 12 个内容项中的 9 个标注了十八个高价值字段，并排除了普通叙述文本和偶然出现的章节编号。

\subsection{攻击者能力}

攻击能力并不构成严格的单调层级，而是由曝光、拍摄几何条件、时间采样和后处理能力共同决定。本文按以下范围报告结果：
\begin{itemize}
\item \textbf{主要范围}：在标称 3.91\,ms UVC 曝光控制设置下进行仪器化快照采集，随后使用 Surya、EasyOCR 和 Tesseract 逐帧识别。主要结果使用原始相机裁剪图像。另一个三引擎压力测试允许攻击者针对每次采集，在五种预先声明的确定性变换和原始输入中进行选择。该设置模拟的是固定 UVC 设置下的孤立图像和有界预处理网格，而非普通的自动曝光手机快照、覆盖完整显示周期的观测，或能够调整无约束增强管线的针对性操作者。实验未记录环境照度。这是一个狭窄的实测条件，并不代表能力最强的合理攻击者。
\item \textbf{实际失效边界}：长曝光、连续录像后的时域视频平均、离轴拍摄、已知周期下的相位搜索与通道重构、目标检测/关联任务，以及商业视觉语言模型的端到端识别（\S\ref{sec:vlm}）。所评估的视频攻击对 150 帧短曝光图像取平均，按 60\,fps 计算名义时长为 2.5\,s。它证明持续录像可以构成一种实际绕过方式，但没有证明可从任意四帧突发序列中恢复信息。本文量化这些风险，但仅将高抑制档结果解释为缓解效果，而非阻断保证。
\item \textbf{尚未得到充分覆盖}：其他相机/显示器组合、精确的滚动快门行对齐、真实场景中的时域超分辨率、以显示序列训练的学习式重构器，以及长期自适应攻击。
\end{itemize}

\subsection{攻击者知识}

攻击者\emph{知晓}时序掩蔽算法、子帧数~$n$ 以及候选周期参数，但\emph{不知道}当前周期的掩码种子或相位起点。密码学安全伪随机数生成器（CSPRNG）为每个周期提供掩码随机化，以防止固定模式学习和跨周期相关；但它不对全周期平均攻击提供密码学保证。一旦攻击者获取并配准一个完整周期，未知种子便无法阻止线性重构。稀疏采样可能破坏传统 OCR 所依赖的字形证据，但当前设计并未将二值化、分割或识别中的任何一个阶段确定为责任环节。端到端 VLM 从同一批归档片段中恢复出了大量文本（\S\ref{sec:vlm}），因此传统 OCR 观察结果不能迁移至这些模型。

\subsection{保护目标}

本文采用短曝光下的分层设计目标来解释所评估的档位。高抑制档以传统 OCR 字符恢复率 $<5\%$ 且完全匹配率 $=0\%$ 为目标；可读性优先档则以相对降低至少 80\% 且完全匹配率 $<1\%$ 为目标。这些论文层面的阈值并非预注册的安全保证。它们不适用于其他曝光条件、自适应预处理、VLM 或其他设备。目标检测在交并比为 0.50 时的平均精度均值（mAP50）变化属于补充诊断，而非通过标准。$\Delta E_{00}$ 和结构相似性（SSIM）数值是数字重构代理，并非真实面板感知或用户舒适度的证据。


%% ============================================================
%% IV. 系统设计
%% ============================================================
\section{系统设计}
\label{sec:method}

\subsection{时序子帧分解}

给定归一化原始帧 $\mathbf{I}\in[0,1]^{H\times W\times C}$，我们生成二值掩模 $\mathbf{M}_1,\ldots,\mathbf{M}_n$，使其满足 $\sum_{k=1}^{n}\mathbf{M}_k=\mathbf{1}$，并令 $\mathbf{I}_k=\mathbf{I}\odot\mathbf{M}_k$。由此得到：
\begin{itemize}
\item \textbf{数字域完备性}：$\sum_{k=1}^{n}\mathbf{I}_k = \mathbf{I}$；
\item \textbf{互斥性}：每个像素恰好在 $n$ 个时隙中的一个时隙内被点亮。
\end{itemize}
这里必须区分“求和”与“时域平均”。对于时隙持续时间和峰值亮度均相同的理想线性显示器，一个周期内的平均辐亮度是 $\mathbf{I}/n$，而不是 $\mathbf{I}$。数字域完备性表示可以对各子帧重新求和；它并不意味着真实显示器的亮度得到无损保留。

这一区别同时构成实验局限。物理拍摄研究未纳入亮度匹配的静态调暗对照、与增强档相匹配的无噪声版本，也未进行面板时序的光度测量。因而，观察到的档位级下降不能仅归因于时序稀疏性；占空比亮度、静态条纹/字形叠加、面板响应、相机曝光以及图像信号处理器（ISP）处理仍相互混杂。

像素分配由基于 ChaCha20~\cite{b_chacha2008} 的密码学安全伪随机数生成器（CSPRNG）产生。拒绝采样将字节流映射为无偏的时隙索引；另行采用 Fisher--Yates 洗牌~\cite{b_knuth1997} 随机化子帧播放顺序。卡方时隙计数检验（$\alpha=0.01$，$\mathrm{df}=n{-}1=3$）仅用作实现的健全性检查，最多重新采样 5 次，并不构成密码学强度的证据。随机化可减少固定图案与跨周期相关性，但无法阻止对已配准完整周期的线性重建。图~\ref{fig:pipeline} 展示了完整管线。

\begin{figure*}[!t]
\centering
\includegraphics[width=\textwidth,viewport=0 0 515.52 206.16,clip]{figures/visio/figure2_method_pipeline/final/figure2_method_pipeline.pdf}
\caption{从源内容到受保护显示的时序像素掩蔽系统管线。输入帧被分解为由 CSPRNG 分配的互补像素掩模；随后在高刷新率播放前，插入可选的面向 OCR 的互补噪声、与档位相关的条纹/字形叠加，以及可选的部分反色帧。人类观察者预期可将多个时隙整合为可读视图，而短曝光相机可能仅捕获一个稀疏子帧。长曝光和视频时域平均能够累积多个时隙，因此本文将其作为边界攻击单独评估，而非假定其不存在。}
\label{fig:pipeline}
\end{figure*}

\subsection{互补对抗噪声}

我们在子帧上叠加对抗噪声 $\mathbf{N}_1, \ldots, \mathbf{N}_n$，使其满足：
\begin{equation}
\sum_{k=1}^{n}\mathbf{N}_k = \mathbf{0}
\label{eq:noise_complement}
\end{equation}
式~(\ref{eq:noise_complement}) 仅在求和、裁剪和显示传递函数之前的理想线性域中严格成立。已有研究表明，较小的文本图像扰动即可显著误导 OCR 系统~\cite{b_song2018}，这为面向 OCR 的噪声设计提供了依据。噪声方向由快速梯度符号法（FGSM）式~\cite{b_goodfellow2015} 代理梯度生成，幅度上界为 $\varepsilon=8/255$，随后在时域中进行互补分解。由于显示器无法发出负辐亮度，实现中采用了基底电平和数值裁剪。因此，在伽马校正、量化、裁剪和相机成像之后，噪声不保证能在物理上相互抵消。

\textbf{组件作用}：随机点阵掩模是主要时序机制。互补噪声是理想域中的可选组件；条纹/字形叠加与部分反色是下述增强档使用的非线性附加项。本文保留归档复合档位的定义，不在事后移除其中的组件。这样可以保持与已拍摄数据的对应关系，但不能证明各组件的收益或最优性。在物理拍摄结果中，噪声效果并不单调，且未拍摄与增强档相匹配的无噪声版本。因此，我们解释的是完整档位，而不建议将互补噪声作为默认组件。

\subsection{亮度补偿}

每个像素在一个 $n$ 时隙周期内仅被点亮一次。如果显示器能够在像素激活的时隙内将线性光域的光输出提高到 $n$ 倍，则周期平均亮度将接近原始亮度。当前概念验证（PoC）既未实现逐时隙动态背光协同，也未进行光度验证。所有亮度对齐、CIEDE2000 $\Delta E_{00}$~\cite{b_ciede2000} 和 SSIM~\cite{b_wang2004} 结果均来自数字积分/归一化模型，仅用于跨配置比较，不构成真实面板在光学上等效或用户感到舒适的证据。

\subsection{时序与带宽}

对于仅包含 $n$ 个互补子帧的基础周期，若以 60\,Hz 作为最低周期率，则要求 $f_r\geq60n$；当 $n=4$ 时，对应 240\,Hz。如果周期还包含一个反色帧，则要求变为 $f_r\geq60(n+1)$，在 240\,Hz 下实际周期率将降至仅 48\,Hz。因此，240\,Hz 满足基础四子帧配置，却不能满足含反色帧配置的 60\,Hz 周期目标。人类闪烁感知还受亮度、对比度、视网膜偏心度、空间内容和眼动影响~\cite{b_cai2024,b_davis2015}；仅凭标称周期率无法预测舒适度。

观察者能否将互补子帧整合成感知完整的图像，取决于时间对比敏感函数（TCSF）。临界闪烁融合频率（CFF）并非普适的截止值：Davis 等人~\cite{b_davis2015} 发现，对于均匀调制，人眼在常规频率附近仍具有敏感性；而在存在高频空间边缘时，超过 500\,Hz 仍可能被观察到。因此，含反色配置的 48\,Hz 完整周期率可能令部分观察者察觉闪烁，基础配置的 60\,Hz 也不能被假定为无闪烁。只有计划中的用户研究才能确定这两种配置是否支持舒适的时序整合。

\subsection{探索性高抑制档}
\label{sec:hardened}

基础线性时序掩模在完整周期时域平均下会发生泄漏（见 \S\ref{sec:security_frontier}）。为应对此问题，我们引入非线性增强档。除非另有说明，本小节各档位均采用 $n=4$ 的基础时序掩蔽，并启用互补对抗噪声。表~\ref{tab:profile_composition} 汇总全部档位配置；表~\ref{tab:real_ocr_full_pool} 中的“掩蔽加噪声”行是不含条纹/字形抗 OCR 叠加的基础条件。

\begin{table*}[!t]
\caption{档位组成汇总（所有档位均取 $n=4$）}
\label{tab:profile_composition}
\centering
\footnotesize
\setlength{\tabcolsep}{3.2pt}
\begin{tabular}{lccccccc}
\toprule
档位 & 掩模 & 噪声 & 单元（px） & 条纹宽度（px） & 条纹 $\alpha_s$ & 字形 $\alpha_g$ & 反色 $\alpha$ \\
\midrule
仅掩蔽 & \checkmark & -- & 1 & -- & -- & -- & -- \\
掩蔽加噪声 & \checkmark & \checkmark & 1 & -- & -- & -- & -- \\
Strong & \checkmark & \checkmark & 1 & 10 & 0.10 & 0.12 & -- \\
可读性优先档 & \checkmark & \checkmark & 1 & 10 & 0.10 & 0.12 & 0.2 \\
高抑制档 & \checkmark & \checkmark & 2 & 6 & 0.42 & 0.55 & 0.2 \\
\bottomrule
\end{tabular}
\end{table*}

\begin{itemize}
\item \textbf{Strong（抗 OCR）档}：基础掩模加噪声，掩模单元大小为 1~像素，条纹宽度为 10~像素，条纹幅度 $\alpha_s=0.10$，字形幅度 $\alpha_g=0.12$，不含反色帧。

\item \textbf{可读性优先档}：在 Strong 档基础上加入一个 $\alpha=0.2$ 的反色帧，本文将其指定为面向可用性的复合配置。归档结果表与拍摄元数据使用标签“deployed”，但这一旧标识既不表示部署建议，也不是经过实证优化的证据。

\item \textbf{高抑制档}：一种探索性的安全导向压力测试档，采用基础掩模加噪声，将掩模单元大小增至 2~像素，条纹宽度设为 6~像素，条纹和字形幅度分别为 0.42 和 0.55，并加入一个 $\alpha=0.2$ 的反色帧。
\end{itemize}

非线性档偏离严格完备性（$\sum\mathbf{I}_k \neq \mathbf{I}$），会产生可见条纹与颗粒。这一取舍与 Kaleido 类时序重编码相反：后者依赖周期内在感知上完整的混合来维持观看体验，这意味着完整周期积分必然重建出原始画面的近似结果，与其防盗录目标相符。而在本文处理的静态内容拍摄威胁中，完备性恰恰是积分攻击的切入点；高抑制档主动将破坏严格完备性作为设计取舍，以可见条纹和颗粒为代价，换取完整周期积分下较低的恢复率。其安全收益和未经测量的可读性代价必须分别报告，也不应将其视为经过用户验证的候选部署方案。

本文采用“高抑制”这一论文表述，因为它只描述该档位在表~\ref{tab:profile_composition} 中的相对位置，并不暗示已经达到某项抗拍摄目标，也不暗示用户能够容忍该档位。为保证可复现性，归档标识仍为 \texttt{capture\_hardened}（旧标签为 \texttt{vlm}）；\S\ref{sec:real_capture} 报告其未达到目标以及残余泄漏。

\subsection{部分反色帧}

当曝光窗口覆盖一个完整显示周期时，即构成长曝光攻击（$n=4$、240\,Hz 时，基础周期约为 16.7\,ms；加入反色帧后，$n+1$ 时隙周期约为 20.8\,ms）。归档中的长曝光标签 31.25 和 125\,ms 是根据 UVC 控制量推导的标称值，而非光度测量值。在归一化域 $\mathbf{I}\in[0,1]$ 中，实现会在每个 $n$ 子帧周期后插入一个\textbf{部分反色帧} $\alpha(\mathbf{1}-\mathbf{I})$。其等价的 8 位表达式为 $\alpha(255-\mathbf{I}_{8})$。增大 $\alpha$ 会同时改变积分信号，以及数字重建结果的亮度和对比度。

真机拍摄的反色强度消融（表~\ref{tab:inversion_ablation}；管线见 \S\ref{sec:real_capture}）显示，$\alpha=0$ 时恢复率为 68.6\%，$\alpha=0.2$ 时为 67.0\%。在 $\alpha\in[0,0.5]$ 范围内，各点估计值介于 61.7\% 和 72.0\% 之间；拍摄样本重采样区间仅用于描述，区间重叠也不是等效性检验。该消融固定 Strong 叠加层，并使用由五项账号/代码/CET6 内容组成的子集。其 $\alpha=0.2$ 数值（$N=135$）并非对可读性优先行的重复估计。接近全反色（$\alpha=1.0$）时，观察到的恢复率降至 18.1\%，但数字积分视图也严重退化。在 150 帧时域均值下，弱反色设置仍保留 64--73\% 的恢复率。可读性优先档选择 $\alpha=0.2$ 是一种模型化设计决策，而非经过验证的最优值。

\begin{table}[!t]
\caption{固定 Strong 叠加层（$\alpha_s=0.10$，$\alpha_g=0.12$）上的反色帧强度 $\alpha$ 消融：真机拍摄引擎最优字符恢复率（均值与 95\% 拍摄样本重采样区间，汇总 9 个几何条件下的五项账号/代码/CET6 内容；长曝光每个设置 $N=135$，视频时域平均每个设置 $N=45$）。该五项子集不同于表~\ref{tab:real_ocr_full_pool} 中可读性优先行使用的 12 项内容。}
\label{tab:inversion_ablation}
\centering
\footnotesize
\begin{tabular}{lcc}
\toprule
$\alpha$ & 长曝光字符（\%） & 视频均值字符（\%） \\
\midrule
0.0 & 68.6~[63.1, 73.9] & 72.7~[63.3, 81.1] \\
0.2 & 67.0~[61.3, 72.7] & 69.8~[59.9, 78.7] \\
0.3 & 72.0~[66.2, 76.9] & 64.3~[54.0, 73.6] \\
0.5 & 61.7~[55.6, 67.7] & 66.4~[56.5, 75.2] \\
1.0 & 18.1~[14.0, 22.0] & 28.4~[20.5, 36.9] \\
\bottomrule
\end{tabular}
\end{table}

因此，每个显示周期包含 $n+1$ 个时隙。当 $n=4$ 且刷新率为 240\,Hz 时，周期率为 48\,Hz，低于 60\,Hz 的设计目标，并可能产生可感知闪烁；不存在适用于所有观看条件的固定“无闪烁阈值”。

\subsection{实现范围}

当前 PoC 工作在应用/算法层，使用 Python 和图形处理器（GPU）渲染来生成显示序列。它不包含操作系统（OS）驱动级注入（DXGI / KMS-DRM / CoreDisplay），也不包含逐时隙动态背光控制。本文验证的是显示序列及其拍摄结果，并不声称已实现生产级系统集成。


%% ============================================================
%% V. 实验评估
%% ============================================================
\section{实验评估}
\label{sec:evaluation}

\subsection{实验设置}
\label{sec:setup}

\textbf{硬件平台}：大多数实验使用一台配备 Intel Core i7-8700K（3.70\,GHz）、32\,GB RAM 和 NVIDIA TITAN Xp（12\,GB VRAM）的桌面工作站。归档的 d0.5/$0^{\circ}$ 运行是在另一台采集机器上使用 $2560\times1600$ 播放画布完成的，但该机器的硬件规格未被保留；其余运行均使用 $1920\times1080$。显示器为 AOC 24G51Z，以标称 240\,Hz 驱动。标称 240-Hz 下，每个时隙持续 4.17\,ms；基础周期和加入反色帧后的周期率分别为 60 和 48\,Hz。面板转换行为未被测量，也不能根据刷新率规格推断。

\textbf{拍摄设备}：所有物理拍摄均使用一台 eMeet SmartCam S600 USB 网络摄像头及其标准 UVC 驱动。相机帧在识别前经过透视校正和内容裁剪。逐次拍摄元数据保存了由 DirectShow/UVC log$_2$ 控制量推导的曝光标签：$-11$ 和 $-8$ 分别标称映射为 0.49 和 3.91\,ms，$-5$ 和 $-3$ 分别映射为 31.25 和 125\,ms。八个几何条件使用共同的 $-8/-5$ 设置；d0.5/$15^{\circ}$ 使用 $-11/-3$。这些数值是推导得到的控制标签，并非实测快门时间。采集代码尝试采用手动曝光，但归档元数据未保存经过语义解码的自动曝光状态。增益、自动白平衡和对焦状态也未被保留。整个采集过程中，显示器亮度设置与采集相位保持固定，但未测量显示器的物理亮度。视频以标称 60\,fps 拍摄。主要 OCR 输入为经过几何校正的内容裁剪区域。

相机与显示器的距离分别为 0.5\,m、1\,m 和 1.5\,m，并在每个距离下分别旋转至 $0^{\circ}$、$15^{\circ}$ 和 $30^{\circ}$。环境照度未被记录。厂商标明该设备采用 8 百万像素 4K 传感器，并支持 1080p/60\,fps 输出~\cite{b_emeet_s600}，但没有说明传感器型号。滚动快门读出、面板输出和应用层时隙时序均未经过光度表征。行混合与面板转换可能解释仿真到真机之间的差距，但现有数据无法识别各自的贡献。结果仅适用于这一 UVC 链路。

\textbf{测试内容}：每个位置均使用 12 个固定内容项：两个账号字符串、两个代码片段、两个数字字符串、一个英文句子、两个混合语言片段、两个英文段落，以及一份 CET6（大学英语六级考试，College English Test Band~6）文档。CET6 页面含有中文文本，两个混合语言内容项也含有中文字符；因此，该语料并非全英文或全 ASCII。11 个合成内容项旨在覆盖上述内容形式，但从 120 个样本的仿真语料中选择这些内容项的过程并未预注册。三次重复使每个档位--攻击--位置单元包含 36 次拍摄；计入已披露的重复采集之前，9 个位置合计为 324 次。检测与跟踪使用不同的子集，仅作为补充诊断。

\textbf{评估指标}：OCR 字符恢复率定义为 $\max(0,1-\mathrm{CER})$，其中 CER 为字符错误率。完全匹配率在移除空白后，以不区分大小写的全文比较计算。显式字段恢复使用评分前固定的 18 个字段，覆盖 9 个含字段内容项。仅当某字段完整的 Unicode 归一化形式出现在该次拍摄至少一个引擎的输出中时，才将其计为已恢复；比较时忽略大小写、空白和标点。本文分别报告字段级微平均恢复率和样本级宏平均恢复率。泄漏率是字符恢复率 $\geq20\%$ 的样本比例，仅作为描述性阈值。在匹配的共同设置单元上，可读性优先档字符恢复率的 P50、P75、P90、P95 和 P99 分别为 10.1\%、22.0\%、45.6\%、63.6\% 和 95.0\%。高抑制档对应值分别为 3.4\%、9.2\%、15.2\%、19.1\% 和 22.6\%。

\textbf{OCR 聚合与不确定性}：对于每次拍摄，字符指标与全文指标分别在三个引擎中取最大值。显式字段采用三个引擎所恢复字段的并集。在匹配前，先对档位、内容、位置和重复编号相同的重复拍摄取平均。主要比较使用三个档位中都存在的键。配对对比将 12 个内容项作为聚类进行重采样，共 2,000 次，随机种子为 20260612。这些内容项聚类区间将所测几何条件视为固定，不能估计新设备上的不确定性。表~\ref{tab:real_ocr_full_pool} 则给出不平衡归档的拍摄样本重采样汇总；这些区间用于描述，而非总体置信区间。在所有几何条件上，未保护条件的恢复率比可读性优先档高 78.0 个百分点（[75.5, 80.3]），比高抑制档高 89.1 个百分点（[85.0, 92.4]）。仅掩蔽档比可读性优先档高 3.0 个百分点（[-0.5, 6.5]），因此归档数据不能证明每个组件都具有独立收益。

\textbf{共同设置分析}：d0.5/$15^{\circ}$ 位置采用了不同的标称 UVC 设置，并产生了明显欠曝的受保护帧。其余 8 个几何条件构成主要固定设置结果（表~\ref{tab:real_ocr_common}）。经过重复采集平均和三个档位匹配后，每个档位均有 $N=288$ 个单元。未保护、可读性优先和高抑制档的平均恢复率分别为 94.5\%、17.8\% 和 5.6\%。配对对比分别为 76.7 个百分点（[74.3, 79.2]）和 88.9 个百分点（[85.5, 91.9]）。可读性优先档在不平衡拍摄图像集合中的平均值为 16.7\%，共 408 张图像，仅作为敏感性汇总报告。排除不可比较的几何条件后，可读性优先档的 150 帧视频均值恢复率由 71.1\% 升至 79.7\%，高抑制档则由 47.9\% 升至 53.9\%。

\textbf{其他稳健性检查}：对五个重复内容项进行留一轮分析，可读性优先档的短曝光恢复率在第 1 轮为 16.6\%，在第 2 轮为 15.8\%，档位级解释未发生变化。检测与跟踪差异尚未进行显著性检验，仅在补充材料中作描述性报告。

\textbf{固定预处理压力测试}：我们在全部 984 张选定的主要样本池图像上评估了 Tesseract、EasyOCR 和 Surya。固定输入网格包含原始输入、gamma 0.5、亮度对比度受限自适应直方图均衡化（CLAHE；第 1--99 百分位拉伸，裁剪限值 4.0，$8\times8$ 分块）、非锐化掩蔽（$\sigma=1.0$，权重 1.7/$-0.7$）、高斯自适应阈值化（块大小 31，常数 5），以及 2$\times$ 双三次上采样。所有参数均在评分前固定。对于每项指标和每次拍摄，攻击者保留最佳引擎--输入组合；随后先对重复拍摄取平均，再执行与主要分析相同的三个档位匹配。完整矩阵包含 17,712 个单元（984 张图像 $\times$ 6 种输入形式 $\times$ 3 个引擎），且未出现 OCR 错误。非原始输入的 Tesseract 单元通过 \texttt{pytesseract 0.3.13} 调用 Tesseract 5.4.1；EasyOCR 1.7.2 和 Surya 0.14.7 的变换输入单元在上述工作站上使用 CUDA。原始输入单元从主要归档导入，其中未保存 Windows 版 Tesseract/EasyOCR 的确切构建版本。该固定网格最优选择（oracle）在声明的网格内对攻击者有利，但它不是无约束自适应搜索、学习式增强攻击或商业 OCR 上界。

数字重建代理指标：$\Delta E_{00}$ 是原始图像与完整周期数字重建图像之间的平均 CIEDE2000 色差；SSIM~\cite{b_wang2004} 是二者的平均结构相似度。两项指标均在数字管线中计算，仅用于本文内部的档位比较。两者均不能衡量真实面板的光学输出、时序舒适度或人类可读性。

% ----------------------------------------------------------
\subsection{真机拍摄 OCR 实验}
\label{sec:real_capture}

物理拍摄归档共包含 10,575 张图像（每个位置 1,175 张），全部来自 S600 网络摄像头。该总数混合了消融、参数扫描与主要攻击条件。可读性优先档对 12 个内容项中的 5 项包含第二轮采集，因此全位置短曝光/视频图像数为 459/153，而不是 324/108。匹配对比前先对重复键取平均。没有采集第二批长曝光数据。Surya~\cite{b_surya2024}、EasyOCR~\cite{b_easyocr2020} 和 Tesseract~\cite{b_smith2007} 共生成 31,725 条引擎级记录。Surya 重跑固定使用 \texttt{surya-ocr 0.14.7}；结果归档未嵌入 Windows 版 Tesseract 和 EasyOCR 的确切构建版本，这仍是一项可复现性局限。每个几何条件使用归档的四点单应变换和输出画布，随后应用已记录的内容裁剪；主要 OCR 管线不包含其他自适应增强。各档位按照 original、mask-only、mask+noise、Strong、可读性优先和高抑制的固定顺序采集，没有随机化。时间漂移或批次漂移可能混杂档位比较。主要比较采用 8 个共同设置几何条件，而表~\ref{tab:real_ocr_full_pool} 保留全部 9 个条件作为敏感性证据。图~\ref{fig:montage} 展示了 1.0\,m/$0^{\circ}$ 下的数字重建结果和匹配的共同设置拍摄结果。

\begin{table*}[!t]
\caption{共同标称 3.91-ms UVC 设置下主要匹配短曝光 OCR 比较。匹配前先对重复拍摄取平均；每个档位的 $N=288$ 由 12 个内容项 $\times$ 8 个几何条件 $\times$ 3 次重复组成。排除不可比较的 d0.5/$15^{\circ}$ 运行。}
\label{tab:real_ocr_common}
\centering
\footnotesize
\setlength{\tabcolsep}{4.6pt}
\begin{tabular}{lcccl}
\toprule
档位 & $N$ & 平均字符恢复率 & 未保护 $-$ 该档位 & 解释 \\
\midrule
原始（未保护） & 288 & 94.5\% & -- & 共同设置基线 \\
可读性优先档 & 288 & 17.8\% & 76.7\,个百分点 [74.3, 79.2] & 传统 OCR 恢复下降，但上尾仍有失效 \\
高抑制档 & 288 & 5.6\% & 88.9\,个百分点 [85.5, 91.9] & 探索性压力测试档；未达到严格的 $<5\%$ 目标 \\
\bottomrule
\end{tabular}
\end{table*}

\begin{table}[!t]
\caption{主要匹配短曝光样本池中的显式敏感字段完整恢复率。18 个字段在评分前固定；每个档位的 288 个单元中，有 216 个至少包含一个字段。}
\label{tab:explicit_fields}
\centering
\footnotesize
\begin{tabular}{lcc}
\toprule
档位 & 字段微平均（432 次机会） & 样本宏平均 \\
\midrule
原始（未保护） & 91.4\% & 86.8\% \\
可读性优先档 & 17.7\% & 17.5\% \\
高抑制档 & 1.2\% & 1.1\% \\
\bottomrule
\end{tabular}
\end{table}

\begin{table*}[!t]
\caption{单一 UVC 真机拍摄 OCR 敏感性样本池（引擎最优；$N$ 为每行拍摄数；包含全部 9 个几何条件）。方括号内为 95\% 拍摄样本重采样汇总，而非总体置信区间。视频行对 150 帧取平均。主要匹配估计见表~\ref{tab:real_ocr_common}。}
\label{tab:real_ocr_full_pool}
\centering
\footnotesize
\begin{tabular}{lllccc}
\toprule
档位 & 攻击模式 & $N$ & 字符恢复率 [95\% 汇总] & 完全匹配率 & 泄漏率 \\
\midrule
\multirow{3}{*}{原始（未保护）}
 & 短曝光 & 324 & 94.1 [93.1, 95.0]\% & 65.7\% & 99.4\% \\
 & 视频（150 帧均值） & 108 & 79.9 [73.5, 86.0]\% & 61.1\% & 86.1\% \\
 & 长曝光 & 324 & 47.3 [42.8, 51.9]\% & 18.8\% & 58.6\% \\
\midrule
\multirow{3}{*}{仅掩蔽}
 & 短曝光 & 324 & 19.2 [17.0, 21.4]\% & 0.3\% & 30.6\% \\
 & 视频（150 帧均值） & 108 & 79.0 [73.0, 84.9]\% & 48.1\% & 88.0\% \\
 & 长曝光 & 324 & 67.2 [63.0, 71.5]\% & 35.5\% & 75.0\% \\
\midrule
\multirow{3}{*}{掩蔽加噪声}
 & 短曝光 & 324 & 25.6 [22.5, 28.8]\% & 2.8\% & 36.1\% \\
 & 视频（150 帧均值） & 108 & 79.2 [73.3, 85.0]\% & 49.1\% & 88.9\% \\
 & 长曝光 & 324 & 68.0 [63.7, 72.4]\% & 39.8\% & 76.2\% \\
\midrule
\multirow{3}{*}{Strong（抗 OCR）}
 & 短曝光 & 324 & 20.7 [18.2, 23.3]\% & 0.6\% & 32.4\% \\
 & 视频（150 帧均值） & 108 & 78.6 [72.5, 84.5]\% & 50.9\% & 88.9\% \\
 & 长曝光 & 324 & 61.6 [57.1, 66.1]\% & 39.8\% & 71.3\% \\
\midrule
\multirow{3}{*}{可读性优先档}
 & 短曝光 & 459 & 15.1 [13.3, 17.1]\% & 0.4\% & 20.5\% \\
 & 视频（150 帧均值） & 153 & 71.1 [65.7, 76.2]\% & 42.5\% & 85.0\% \\
 & 长曝光 & 324 & 60.9 [56.6, 65.6]\% & 36.4\% & 69.1\% \\
\midrule
\multirow{3}{*}{高抑制档}
 & 短曝光 & 324 & 5.0 [4.3, 5.8]\% & 0.0\% & 3.7\% \\
 & 视频（150 帧均值） & 108 & 47.9 [41.8, 54.0]\% & 5.6\% & 76.9\% \\
 & 长曝光 & 324 & 9.3 [7.9, 10.9]\% & 0.0\% & 14.2\% \\
\bottomrule
\end{tabular}
\end{table*}

\textbf{主要发现}：

\begin{enumerate}
\item \textbf{主要固定设置证据}：表~\ref{tab:real_ocr_common} 给出主要短曝光估计。每个档位均使用相同的 288 个匹配单元（12 个内容项 $\times$ 8 个几何条件 $\times$ 3 次重复），未保护、可读性优先和高抑制档的平均恢复率分别为 94.5\%、17.8\% 和 5.6\%。配对的内容项聚类对比分别为 76.7 个百分点（[74.3, 79.2]）和 88.9 个百分点（[85.5, 91.9]）。包含全部 9 个几何条件的拍摄图像池给出的数值为 94.1\%、15.1\% 和 5.0\%，但该样本池混合了不同曝光设置和重复采集。因此，得到支持的结果仅限于这一单相机、标称设置管线。

\item \textbf{固定预处理实质性削弱了原始输入结果}：在每个档位相同的 288 个匹配单元上，原始输入的三引擎最优选择给出 94.5\% 的未保护恢复率、17.8\% 的可读性优先档恢复率和 5.6\% 的高抑制档恢复率。允许五种固定变换后，这些数值分别升至 95.9\%、40.2\% 和 13.7\%（表~\ref{tab:preprocessing_grid}）。在这一更强的传统 OCR 条件下，未保护减可读性优先的配对差值为 55.7 个百分点（[50.3, 60.5]），未保护减高抑制的差值为 82.2 个百分点（[78.5, 85.2]）。包含 17,712 个单元的矩阵没有出现 OCR 错误。这些结果排除了“原始输入数值构成传统预处理攻击上界”的主张，但仍未评估无约束增强和商业 OCR。

\item \textbf{档位组件与目标边界}：在汇总归档中，仅掩蔽条件达到 19.2\%，掩蔽加噪声达到 25.6\%，Strong 达到 20.7\%。因此，归档不能证明噪声或弱叠加具有单调收益。可读性优先档是一项复合观察结果，而不是每个组件均有帮助的证据。在匹配样本池中，该档位仍有 17.7\% 的字段微平均完整恢复率（表~\ref{tab:explicit_fields}）和 95.0\% 的 P99 字符恢复率。高抑制档同样未达到其描述性的 $<5\%$ 目标；在共同标称设置下，其匹配平均值为 5.6\%。

\item \textbf{积分攻击仍是残余失效边界}（图~\ref{fig:bar_attack}）：在标称长曝光下，可读性优先档的字符恢复率为 60.9\%，完全匹配率为 36.4\%；在 150 帧均值下，二者分别为 71.1\% 和 42.5\%。排除 d0.5/$15^{\circ}$ 后，视频恢复率升至 79.7\%。弱反色（$\alpha=0.2$）并未产生较大的观测积分恢复下降。仅掩蔽档和 Strong 档在 150 帧均值下也较为接近（79.0\% 对 78.6\%）。这些比较仅用于描述，不能隔离组件因果效应。高抑制档的积分恢复率较低，但仍存在大量泄漏。

\item \textbf{长曝光反转是一项值得注意的负面结果}：未保护条件的长曝光字符恢复率基线仅为 47.3\%，低于可读性优先档的 60.9\%。在唯一采用 125\,ms 的位置（0.5\,m/$15^{\circ}$），未保护恢复率为 1.1\%，可读性优先档则为 46.3\%；全亮图像中观察到的饱和截断支持以下解释：亮度降低减少了饱和，使信号移入传感器近似线性的响应区域。排除该位置后，31.25\,ms 异常仍随位置变化：在 0.5\,m 下，结果符合预期方向（未保护 78--82\% $>$ 可读性优先 25--55\%）；而在六个 1.0--1.5\,m 位置中，有五个位置的可读性优先档恢复率高出 1.9--59\,个百分点（其中四个超过 12\,个百分点）。由于缺少匹配的局部图像诊断，我们仅\emph{假设}残余条纹/字形边缘会与几何条件特定的过曝或对比度压缩相互作用；现有数据不能确立该机制。这一异常不削弱短曝光结论，因为可读性优先档在短曝光下始终更低；但它表明固定长曝光参数会产生非单调效应。主要有效性结论限于单帧短曝光；积分攻击作为安全边界报告。

\item \textbf{高抑制档：积分泄漏降低，但感知代价尚未解决}：这一探索性压力测试档在标称长曝光下得到 9.3\% 字符恢复率 [7.9, 10.9] / 0\% 完全匹配率，在 150 帧均值下得到 47.9\% 字符恢复率 [41.8, 54.0] / 5.6\% 完全匹配率 [1.9, 10.2]。方括号内为拍摄样本重采样汇总。汇总均值掩盖了较大的位置间差异，其中包括一个恢复率为 0\% 的欠曝会话。排除该会话后，视频恢复率升至 53.9\% [48.0, 59.9]，完全匹配率升至 6.3\%。任何部署评估都需要覆盖更广的几何条件、曝光、跨设备和用户可接受性证据。用户研究协议不包含该档位。

\item \textbf{几何趋势}：在 9~个条件（0.5--1.5\,m $\times$ 0--30$^{\circ}$）下，受保护短曝光的恢复率始终低于未保护值。由于实验设计不平衡，且未拟合几何因素的推断模型，我们不主张恢复率随距离或角度具有统计显著的单调鲁棒性。
\end{enumerate}

\begin{figure}[!t]
\centering
\includegraphics[width=\columnwidth]{figures/real_capture_montage.pdf}
\caption{1.0\,m、$0^{\circ}$ 条件下的真机拍摄蒙太奇。在纵向堆叠的每个档位组（未保护、可读性优先和高抑制）内，图像从上至下依次为数字积分重建、3.91\,ms 短曝光拍摄和 31.25\,ms 长曝光拍摄。相机图像以相同方式从归档原始相机 JPEG 中裁剪，未进行几何校正或亮度增益，因此与 OCR 评估使用的几何校正裁剪图不同。数字图像对每个档位均使用相同的 $1920\times1080$ 黑色播放画布和亮度对齐的全周期积分；可读性优先和高抑制重建包含记录的 $\alpha=0.2$ 反色时隙。这些数字视图并非人类可读性或视觉舒适度的直接证据。}
\label{fig:montage}
\end{figure}

\begin{figure}[!t]
\centering
\includegraphics[width=\columnwidth]{figures/real_capture_bar.pdf}
\caption{三个代表性真机拍摄档位（未保护、可读性优先和高抑制）在短曝光、长曝光及 150 帧视频时域均值下的引擎最优 OCR 字符恢复率和完全匹配率，汇总全部 9 个所测距离--角度配置。柱越低，表示恢复出的文本越少。长曝光与视频柱揭示了积分边界：即使可读性优先档的短曝光柱较低，在积分攻击下仍保留大量可恢复内容。可读性优先档的样本数包含已披露的五项内容重复采集，因此与其他档位不同。}
\label{fig:bar_attack}
\end{figure}

\textbf{分引擎结果}：表~\ref{tab:real_ocr_engine} 按引擎拆分汇总的原始输入短曝光结果。不同引擎的未保护恢复率差异很大（Surya 37.1\%、EasyOCR 79.5\%、Tesseract 84.3\%）。Strong 条件下的恢复率范围为 2.1--16.3\%，可读性优先档为 3.2--11.7\%，高抑制档为 1.8--2.0\%。表~\ref{tab:preprocessing_grid} 和图~\ref{fig:real_engine_ocr} 则使用匹配的共同设置样本池，并表明每个引擎都至少受益于一种固定变换。对于可读性优先档，Tesseract 的恢复率由 3.2\% 升至 10.8\%，EasyOCR 由 14.2\% 升至 33.9\%，Surya 由 4.1\% 升至 20.2\%。对于高抑制档，相应变化分别为 2.0\% 至 8.2\%、2.3\% 至 6.5\%，以及 2.0\% 至 6.0\%。按指标选择的网格最优结果在这一有限的开源引擎与变换集合内对攻击者有利；本文未测试 Google Cloud Vision、Azure AI Vision 或百度 OCR 等商业 OCR 应用程序接口（API），这些服务可能取得更高恢复率，因此该结果并非现有 OCR 服务的上界。

\begin{table}[!t]
\caption{真机拍摄短曝光 OCR：分引擎结果（汇总 9 个几何条件）}
\label{tab:real_ocr_engine}
\centering
\footnotesize
\setlength{\tabcolsep}{5.5pt}
\begin{tabular}{llccc}
\toprule
档位 & 引擎 & 字符（\%） & 完全匹配（\%） & N \\
\midrule
\multirow{3}{*}{原始}
 & Surya & 37.1 & 22.8 & 324 \\
 & EasyOCR & 79.5 & 27.8 & 324 \\
 & Tesseract & 84.3 & 47.8 & 324 \\
\midrule
\multirow{3}{*}{仅掩蔽}
 & Surya & 4.2 & 0.3 & 324 \\
 & EasyOCR & 15.6 & 0.0 & 324 \\
 & Tesseract & 3.2 & 0.0 & 324 \\
\midrule
\multirow{3}{*}{掩蔽加噪声}
 & Surya & 11.0 & 2.8 & 324 \\
 & EasyOCR & 18.1 & 0.0 & 324 \\
 & Tesseract & 4.8 & 0.0 & 324 \\
\midrule
\multirow{3}{*}{Strong}
 & Surya & 6.1 & 0.6 & 324 \\
 & EasyOCR & 16.3 & 0.0 & 324 \\
 & Tesseract & 2.1 & 0.0 & 324 \\
\midrule
\multirow{3}{*}{可读性优先档}
 & Surya & 3.2 & 0.4 & 459 \\
 & EasyOCR & 11.7 & 0.0 & 459 \\
 & Tesseract & 3.3 & 0.0 & 459 \\
\midrule
\multirow{3}{*}{高抑制档}
 & Surya & 1.8 & 0.0 & 324 \\
 & EasyOCR & 2.0 & 0.0 & 324 \\
 & Tesseract & 1.9 & 0.0 & 324 \\
\bottomrule
\end{tabular}
\end{table}

\begin{table}[!t]
\caption{固定预处理网格下匹配共同设置的字符恢复率。每个单元格均为原始输入 $\rightarrow$ 逐次拍摄网格最优选择（\%）；重复采集平均和匹配后，每个档位 $N=288$。}
\label{tab:preprocessing_grid}
\centering
\footnotesize
\setlength{\tabcolsep}{4.2pt}
\begin{tabular}{lccc}
\toprule
引擎 & 原始 & 可读性优先 & 高抑制 \\
\midrule
Tesseract & 84.9 $\rightarrow$ 91.8 & 3.2 $\rightarrow$ 10.8 & 2.0 $\rightarrow$ 8.2 \\
EasyOCR & 85.1 $\rightarrow$ 88.4 & 14.2 $\rightarrow$ 33.9 & 2.3 $\rightarrow$ 6.5 \\
Surya & 30.6 $\rightarrow$ 49.3 & 4.1 $\rightarrow$ 20.2 & 2.0 $\rightarrow$ 6.0 \\
三者最优 & 94.5 $\rightarrow$ 95.9 & 17.8 $\rightarrow$ 40.2 & 5.6 $\rightarrow$ 13.7 \\
\bottomrule
\end{tabular}
\end{table}

\begin{figure}[!t]
\centering
\includegraphics[width=\columnwidth]{figures/real_engine_ocr.pdf}
\caption{匹配共同设置样本池中，各 OCR 引擎的原始输入与固定网格最优字符恢复率对比（每个档位 $N=288$）。网格最优选择针对每次拍摄和每项指标，分别从原始输入加五种预先声明的变换中选取最佳结果。连线表示同一引擎内部的增幅，而非不确定性区间。确切数值见表~\ref{tab:preprocessing_grid}。}
\label{fig:real_engine_ocr}
\end{figure}
\subsection{真机拍摄 VLM 探针实验}
\label{sec:vlm}

为检验传统 OCR 的结论能否推广至更强的识别器，我们使用了通过 SiliconFlow 应用程序编程接口（API）提供的三个商业视觉语言模型（VLM）端点：Qwen3-VL-32B-Instruct~\cite{b_qwen3vl2025}、Kimi-K2.6 和 GLM-4.5V~\cite{b_glm45v2025}。本节是一项\emph{边界探针}：若 VLM 能够恢复传统 OCR 无法识别的文本，则本方法的实际适用范围仅限于缓解传统 OCR 对短曝光图像的恢复。

探针实验覆盖三个距离下的 0$^{\circ}$ 正轴拍摄。0.5\,m 实验包含 156 个输入：36 个未保护短曝光基线、各 36 个高抑制短曝光和长曝光输入，以及从 60\,fps 视频生成的四种聚合视图（每种 12 个片段）。1.0\,m 和 1.5\,m 实验各包含 120 个高抑制输入：各 36 个短曝光和长曝光输入，外加四类视频聚合视图（每类 12 个片段）。VLM 输入与表~\ref{tab:real_vlm} 中 OCR 的最优引擎（best-of-engine，BoE）参照使用同一批文件，其标称标签为 3.91/31.25\,ms。所有 VLM 调用均使用相机原始 JPEG。\S\ref{sec:real_capture} 中固定的三引擎网格评估的是传统 OCR，而不是 VLM 输入预处理；先导实验中的 CLAHE 增强恢复出了更多 VLM 文本，因此基于原始 JPEG 的 VLM 结果并非攻击上界。

在 1,188 次模型请求中，1,133 次返回了可解析结果：0.5\,m 处失败 34 次，1.0\,m 处失败 21 次，1.5\,m 处无失败。每个表格单元均报告有效 $N$，且结果以响应可解析为条件。GLM 的失败集中在混合内容、数字和凭据条目上，因此不能将缺失视为随机缺失。为此，我们分别将每次失败调用的恢复率赋为 0 或 100\%，计算逐单元格的下界和上界。对于 GLM 在 0.5\,m 的短曝光单元格，字符恢复率的界限为 61.0--74.9\%，完全匹配率的界限为 44.4--58.3\%。我们不使用条件单元格对模型进行排序；所有单元格的完整界限均随分析一并归档。

我们将无失败的 1.5\,m 实验作为主要的跨模型证据；1.0\,m 实验则作为短曝光距离扫描的补充数据。在 1.0\,m 处，Qwen 和 Kimi 的长曝光完全匹配率分别为 84.4\% 和 74.3\%，三个模型的成功调用在 150 帧视频均值下得到 70.0--83.3\% 的完全匹配率。GLM 在 1.0\,m 处的长曝光完全匹配率仅为 3.3\%，而在 1.5\,m 处为 58.3\%；但依赖内容的失败使我们无法主张结果具有跨距离一致性。为保持紧凑，表~\ref{tab:real_vlm} 在 0.5\,m 处仅报告 150 帧视频均值；1.5\,m 实验则报告全部四种聚合视图，以展示结果对视图选择的敏感性。

所有调用均使用 SiliconFlow 的 OpenAI 兼容 API。记录的提供方标识符为 \texttt{Qwen/\allowbreak Qwen3-\allowbreak VL-\allowbreak 32B-\allowbreak Instruct}、\texttt{Pro/\allowbreak moonshotai/\allowbreak Kimi-\allowbreak K2.6} 和 \texttt{zai-\allowbreak org/\allowbreak GLM-\allowbreak 4.5V}。Qwen3-VL 和 GLM-4.5V 有公开技术报告~\cite{b_qwen3vl2025,b_glm45v2025}。对于 Kimi-K2.6，我们以提供方端点标识符和归档的请求元数据作为可复现性记录，而不对另有文档说明的上游公开版本作出主张。第二个标识符中的 \texttt{Pro/} 前缀为 SiliconFlow 所特有，该端点在我们归档的调用中接受图像输入。解码温度为 0，最大输出为 4,096 个词元。共用提示词要求仅输出可见文本并禁止猜测；模型不会获得真实值。完整提示词、参数和逐次调用输出均已归档。指标遵循 \S\ref{sec:setup}，并在拍摄样本层面汇总。公开的 CET6 文档可能因训练数据记忆而抬高恢复率，因此合成片段与完全匹配率能提供更干净的区分性证据。

\textbf{0.5\,m 会话选择透明性}：归档保留了同一位置的另一个会话（\path{real_capture_vlm_d0.5_a0_rerun.json}，时间戳 012715），其中高抑制短曝光输入几乎全黑，因而三个模型均给出空转录，字符恢复率和完全匹配率均为 0\%。表~\ref{tab:real_vlm} 使用 141107（3.91\,ms）会话，因为它与 OCR 参照使用同一批文件，且其输入并非全黑。这是仅有的两个 0.5\,m 正轴 VLM 会话，因此 Qwen 观测到的短曝光完全匹配结果从欠曝光会话的 0/36 个拍摄样本到表中会话的 28/36 个拍摄样本不等，并非稳定的单会话估计。充足曝光准则是在检查两个会话后才采用的，因此属于事后选择规则；保留两个会话并报告结果范围，可以明确揭示其曝光敏感性。表中会话是偏向攻击者的保守情形，另一个会话则保留作为欠曝光诊断。

\begin{table*}[!t]
\caption{真机拍摄 VLM 探针结果（0$^{\circ}$ 正轴；除原始短曝光基线外，所有行均为高抑制档）。每个 VLM 单元格报告完全匹配计数 / 字符恢复率 / $N$（恢复数 / \% / 有效调用数）；存在失败的单元格以响应可解析为条件。OCR BoE 给出同批次三引擎最优引擎参照（完全匹配率 / 字符恢复率，\%）。时域均值行对 150 个短曝光帧取平均；其他聚合视图在正文中定义。0.5\,m 高抑制短曝光行采用曝光充足且偏向攻击者的 141107 会话；仅有的另一个 0.5\,m 会话存在欠曝光，Qwen 的完全匹配结果为 0/36。}
\label{tab:real_vlm}
\centering
\footnotesize
\setlength{\tabcolsep}{4.2pt}
\begin{tabular}{llcccc}
\toprule
距离 & 条件 & OCR BoE & Qwen3-VL-32B & Kimi-K2.6 & GLM-4.5V \\
\midrule
\multicolumn{6}{l}{\textit{短曝光距离扫描}}\\
\midrule
0.5\,m & 原始短曝光 & 58.3~/~92.4 & 30~/~94.5~/~36 & 30~/~95.6~/~36 & 19~/~92.0~/~26 \\
0.5\,m & 短曝光 & 0.0~/~8.4 & 28~/~91.5~/~36 & 17~/~84.7~/~36 & 16~/~70.8~/~31 \\
1.0\,m & 短曝光 & 0.0~/~3.6 & 20~/~79.1~/~36 & 14~/~68.6~/~36 & 9~/~47.0~/~33 \\
1.5\,m & 短曝光 & 0.0~/~6.5 & 17~/~68.9~/~36 & 12~/~61.5~/~36 & 13~/~70.0~/~36 \\
\midrule
\multicolumn{6}{l}{\textit{关键攻击视图}}\\
\midrule
0.5\,m & 长曝光 & 0.0~/~2.1 & 28~/~85.7~/~34 & 1~/~7.9~/~36 & 0~/~4.9~/~30 \\
0.5\,m & 视频（150 帧均值） & 8.3~/~50.7 & 10~/~94.0~/~12 & 10~/~95.5~/~12 & 7~/~91.8~/~8 \\
1.5\,m & 长曝光 & 0.0~/~13.2 & 33~/~89.3~/~36 & 32~/~89.9~/~36 & 21~/~62.4~/~36 \\
1.5\,m & 视频（单帧最优） & 0.0~/~4.3 & 1~/~24.8~/~12 & 0~/~12.7~/~12 & 0~/~29.8~/~12 \\
1.5\,m & 视频（150 帧均值） & 0.0~/~48.4 & 8~/~90.3~/~12 & 8~/~90.3~/~12 & 7~/~88.5~/~12 \\
1.5\,m & 视频（窗口均值最优） & 0.0~/~8.6 & 7~/~89.0~/~12 & 5~/~84.2~/~12 & 8~/~87.2~/~12 \\
1.5\,m & 视频（最大值投影） & 0.0~/~19.0 & 8~/~90.1~/~12 & 7~/~89.2~/~12 & 7~/~88.7~/~12 \\
\bottomrule
\end{tabular}
\end{table*}

\begin{table}[!t]
\caption{0.5\,m 高抑制短曝光条件下按内容类型划分的 VLM 字符恢复率。基线列给出同一会话中 Qwen 对未保护原始短曝光图像的结果（Kimi/GLM 与之相近）；GLM 偶有调用失败，括号内的 $N$ 表示有效数量。在 1.5\,m 处，三个模型对 CET6 密集页面的字符恢复率均为 0\%。}
\label{tab:vlm_content}
\centering
\footnotesize
\setlength{\tabcolsep}{2.8pt}
\begin{tabular}{lcccc}
\toprule
内容类型（$N$） & 基线（\%） & Qwen（\%） & Kimi（\%） & GLM（\%） \\
\midrule
CET6 文档（3） & 64.4 & 18.5 & 5.6 & 0.4 \\
凭据（6） & 97.6 & 97.6 & 88.3 & 77.1~(5) \\
代码片段（6） & 89.8 & 94.9 & 95.8 & 88.9 \\
数字字符串（6） & 100.0 & 100.0 & 98.6 & 50.0~(4) \\
混合内容（6） & 100.0 & 100.0 & 92.5 & 98.2~(4) \\
段落（6） & 100.0 & 99.8 & 81.7 & 66.5 \\
英文句子（3） & 95.3 & 95.3 & 96.9 & 94.6 \\
\bottomrule
\end{tabular}
\end{table}

表~\ref{tab:real_vlm} 给出了汇总的 VLM 探针结果：短曝光行展示距离变化，长曝光和视频聚合行展示积分攻击边界。表~\ref{tab:vlm_content} 按内容类型细分了 0.5\,m 高抑制短曝光条件下的字符恢复率。

\textbf{主要发现}：

\begin{enumerate}
\item \textbf{VLM 限定而非扩展了主张范围}：对于同一批 0.5\,m 照片，在高抑制短曝光条件下，最优引擎 OCR 的字符恢复率仅为 8.4\%，完全匹配率为 0\%；而 Qwen3-VL 的字符恢复率达到 91.5\%，并在来自 12 个内容项的 36 个拍摄样本中完全恢复 28 个。在 1.5\,m 处，三个模型仍保持 33.3--47.2\% 的完全匹配率和 61.5--70.0\% 的字符恢复率。即使提示词禁止猜测，VLM 仍可从可见片段中恢复大量文本。因此，在这项正轴、高抑制、三距离探针实验中，针对传统 OCR 的主张不能推广至 VLM 攻击者；本实验未评估离轴 VLM 拍摄和其他保护档位。

\item \textbf{VLM 攻击下的保护失效集中在大字号短片段}：在 0.5\,m 处，VLM 对 CET6 文档页面的字符恢复率为 0.4--18.5\%，而对凭据、代码和数字字符串的字符恢复率为 50--100\%（表~\ref{tab:vlm_content}）。在 1.5\,m 处，三个模型对 CET6 页面的恢复率均降至 0\%，但数字字符串和短句仍可恢复。这与字体尺度和语言先验的影响一致~\cite{b_shi2023,b_jiang2025}，但本实验并未单独识别其中任何一种机制。

\item \textbf{攻击能力随模型和条件而变化}：在无失败的 1.5\,m 会话中，标称长曝光完全匹配率为 58.3\%--91.7\%，字符恢复率为 62.4--89.9\%。这一差异仅为描述性结果，不应解释为稳定的模型排名。VLM 行为取决于端点版本、距离、拍摄视图和内容。

\item \textbf{持续视频聚合构成 VLM 失效边界}：在 0.5\,m 处，三个模型的 150 帧均值缺失界限覆盖 58.3--91.7\% 的完全匹配率和 61.2--95.5\% 的字符恢复率。所有聚合视图均来自同一段连拍，不能代表相互独立的短时攻击。在无失败的 1.5\,m 会话中，150 帧均值产生 58.3--66.7\% 的完全匹配率和 88.5--90.3\% 的字符恢复率，而单帧最优视图的完全匹配率为 0--8.3\%。本实验未确定所需的最短录制时长。
\end{enumerate}

因此，有效性主张必须由内容和攻击者类型共同限定：在本探针实验中，高抑制档降低了传统 OCR 对密集小字号文本的恢复率；但当攻击者采用现代 VLM 时，大字号短片段（凭据、代码、数字字符串）仍可恢复。这一负面结果将方法的实际适用范围从“阻止任意机器读取屏幕”重新限定为“降低传统 OCR 对短曝光批量图像的恢复”，并将面向 VLM 的显示保护明确为开放问题。对策方向见 \S\ref{sec:vlm_defense}。本节仅覆盖单摄像头、0$^{\circ}$ 正轴拍摄、一个档位、三个商业 API 模型，且每个单元格仅有 12 个视频片段；这些数值是描述性的边界探针结果，而非完整的跨几何条件或跨模型版本评估。

% Defined here for source locality with the real-capture protocol, but emitted
% after all attack experiments so the evaluation narrative remains continuous.
\long\def\deferreduserstudy{
% ----------------------------------------------------------
\subsection{用户体验研究}
\label{sec:user_study}

为评估受保护显示对人类可读性和视觉舒适度的影响，我们设计了一项基于 Web 的用户研究。

\textbf{设备与环境协议}：采用与 \S\ref{sec:setup} 相同的 240\,Hz 显示器。正式研究版本将实测刷新率的硬性门槛设为 $\geq$200\,Hz，以确保固定 $n=4$ 时基本子帧循环频率满足 $f_r/n\geq50$\,Hz；刷新率低于该门槛时将阻止参与，而不是在不提示的情况下减小 $n$。144--199\,Hz 的自适应路径仅用于明确标注的演示模式，其会话默认不纳入分析。每位参与者开始实验前，实验人员确认显示器亮度固定，自动亮度/省电模式和可变刷新率已禁用，浏览器统一为全屏模式，高负载后台进程已关闭，观看距离约为 60\,cm。显示模式发生任何变化后，均重新测量刷新率。

浏览器在每次试验的 \path{requestAnimationFrame} 循环中记录实际帧间隔；若相邻回调的间隔超过根据试验前实测刷新率计算的预期间隔的 1.5$\times$，则记录一次丢帧事件，同时记录估计丢帧数、丢帧率、平均/最大帧间隔、观测刷新率、基本子帧循环频率，以及包含反色帧的全周期频率。这些记录仅反映浏览器的呈现节奏，不能替代对面板扫描或光度学特性的测量。

\textbf{伦理与安全}：拟开展的研究属于最小风险的人体参与研究，目前尚未收集任何参与者数据。招募开始前，作者将从具备权限的机构或院系主管部门取得并记录书面批准、豁免或免审认定；在取得该认定前，不会开始正式数据采集。参与者必须分别勾选知情同意与光敏感筛查：页面明确说明掩蔽显示涉及快速时序调制，若出现不适、眼疲劳、眩晕、恶心或头痛，应立即停止，并可随时无条件退出。自述有光敏性癫痫史或对闪烁敏感者不得继续参与。系统记录同意时间戳和筛查通过标记。实现采用保守的 $f_r/n\geq50$\,Hz 运行下限，但这一工程门槛不构成伦理或临床安全认定。在 240\,Hz 面板上取 $n=4$ 时，基本循环频率为 60\,Hz；反色帧将全周期频率降至 $f_r/(n+1)=48$\,Hz，实现会单独记录警告和实际时序。注册信息（学号、姓名）仅用于参与管理；分析和发表仅使用去标识化数据。年龄和性别为可选的人口统计字段。

\textbf{任务设计}：每位参与者首先完成一次 12 秒、不计分、无掩蔽的 Canvas 打字热身；随后观看 10 秒、不计分的可读性优先完整掩蔽预览；再在同一可读性优先掩蔽下完成一次 8 秒、不计分的打字练习。这一流程将初始刺激适应和输入适应与随后四次各 20 秒的计分试验（被试内设计）分离开来。在每个打字页面上，包括热身、练习和计分试验，参与者点击“开始”前，刺激 Canvas 已可见且播放器已启动。点击“开始”后触发 5 秒倒计时，此时输入仍被禁用；仅在倒计时结束后，计时打字阶段和输入功能才同时开始。控制条件和掩蔽条件均采用相同的开始前观看流程。

两个条件各重复两次。实验人员连续分配的从零开始的序列号 $k$ 用于选择 ABBA 或 BAAB 顺序，以平衡练习效应与疲劳效应：(A)~\textbf{原始文本}（控制条件），以 $n=1$ 的无掩蔽 Canvas 显示；(B)~\textbf{掩蔽文本}，使用可读性优先配置：$n=4$、掩蔽加噪声、强抗 OCR（$\alpha_s=0.10$、$\alpha_g=0.12$、6 个循环）+ 弱反色（$\alpha=0.2$）。打字顺序索引为 $k\bmod2$，主观评分的拉丁方行索引为 $\lfloor k/2\rfloor\bmod6$。因此，每连续 12 名参与者恰好覆盖全部 $2\times6$ 种联合分配。正式会话对 $k$ 强制执行唯一性约束，后端会重新计算并核验两个顺序索引。

两个条件均由同一 \texttt{render\allowbreak Source\allowbreak Canvas} 路径渲染，使用相同的深色背景、浅色文字、Canvas 尺寸、23\,px 粗体字体和颜色。唯一的处理差异是时序保护流水线，从而避免极性和字重混杂。此处的“循环”是指 Web 实现依次轮换六组不同的掩蔽、噪声和条纹实例，随后再重复。这一 Web 播放特有设置与 \S\ref{sec:method} 中的真机拍摄配置共用 $\alpha_s$ 和 $\alpha_g$，并避免反复积分同一固定叠加图案。

评分不再逐位置比较字符，而是使用 Levenshtein 最小字符串距离（minimum string distance，MSD）~\cite{b_soukoreff2003}，将完整转录与最优目标前缀对齐；未输入的目标后缀不计为错误。对齐后得到准确率、每分钟字符数（CPM）和每分钟词数（WPM），同时记录编辑距离、对齐前缀长度和评分版本。首次按键延迟定义为从 5 秒倒计时结束时（此时输入启用且计分计时器启动）到首次使转录内容非空的输入事件之间的时间。分析首先在每名参与者内部对每个条件下的两次试验取平均，随后进行被试内配对比较。

\textbf{主观评分}：打字试验结束后，参与者针对界面显示的三个 1--5 级 Likert 条目，对 6~种条件进行评分：可读性（1 = 难以阅读，5 = 非常清晰）、稳定性（存储在 \texttt{flicker} 字段中；1 = 强烈闪烁，5 = 无可感知闪烁）和即时视觉舒适度（\texttt{fatigue}；1 = 非常不适，5 = 非常舒适）。因此，每个条目的分值越高均代表结果越好。条件包括无掩蔽原始锚点、$n\in\{2,3,4\}$ 的掩蔽加噪声、$n=4$ 的仅掩蔽，以及打字任务所用的同一 $n=4$ 可读性优先完整配置（掩蔽加噪声+强抗 OCR+弱反色）。前三种掩蔽层级在 240\,Hz 面板上均能达到 $\geq$60\,Hz 的真实时序播放；原计划的 $n=8$ 只会产生 30\,Hz 并触发静态回退，因此不纳入时序消融。六种条件按 6 阶平衡拉丁方呈现；每种条件至少观看 10~秒后方可提交，同时记录观看开始/结束时间和持续时长。无掩蔽锚点也用于辅助检查注意力和量表理解，但不会通过“必须评 5 分”的规则机械排除参与者。本研究不包含高抑制档；其可见条纹和颗粒的可接受性留待未来工作评估（见 \S\ref{sec:discussion}）。

\textbf{样本量与分析计划}：计划至少招募 $N=24$ 名参与者，该样本量处于在 80\% 检验效能下检测 $d_z\approx0.6$--$0.8$ 配对效应量所需范围内，并且恰好覆盖两轮每 12 人一次的联合顺序分配。预注册的排除标准包括：调试/演示会话、未完整记录四次打字或六次评分、实测刷新率低于 200\,Hz、任一打字试验尝试字符少于 5 个、控制条件平均准确率低于 50\%、掩蔽试验未处于时序模式、任一试验的观测基本循环频率低于 50\,Hz，或以最短观看时长进行直线式评分。最后一项标准定义为：全部六行评分记录的观看时长均不超过 11~秒（在 10 秒提交门槛上容许 1 秒误差），且每行内部三个条目的分数完全相同。WPM、CPM、准确率、尝试字符数和首次按键延迟首先在每名参与者内部按条件取平均；速度指标必须与准确率联合解释，避免将低准确率下的快速输入误判为性能提升。当预先指定的 Shapiro 检验通过时，使用配对 $t$ 检验检验配对差异；否则使用 Wilcoxon 符号秩检验，并报告效应量和以参与者为单位的 bootstrap 95\% 置信区间。Likert 数据按维度使用 Friedman 检验，并以 Holm 校正的事后成对 Wilcoxon 检验进行分析。分析脚本直接从正式数据库 \texttt{study\_formal.db} 生成去标识化的参与者均值、JSON 审计报告和两个 \LaTeX\ 表格。

% NOTE: All "XX" marks and "--" table cells below are pending data placeholders.
% After the study, run webstudy/analyze_study.py; its analysis_output/typing_table.tex
% and ratings_table.tex rows map one-to-one onto Tables \ref{tab:study_typing} and
% \ref{tab:study_ratings}, and experiments/paper_figures/make_all.py regenerates
% Figs. \ref{fig:study_typing} and \ref{fig:study_ratings} from the same output.
\textbf{参与者}：共招募 XX 名志愿者；其中 XX 个会话满足全部预注册纳入标准并进入分析（排除 XX 个；按标准划分的审计计数随发布的分析输出一并提供）。纳入样本的可选人口统计信息如下：XX 名参与者中有 XX 名报告年龄（中位数 XX 岁，范围 XX--XX），XX 名报告性别（XX 名女性，XX 名男性）。所有正式会话均在 240\,Hz 实验室显示器上运行，实测刷新率至少为 200\,Hz。

\begin{table}[!t]
\caption{打字任务：控制条件与可读性优先掩蔽条件的被试内比较（$N=$ XX；每个条件取每名参与者两次试验的均值；$\Delta$ = 掩蔽 $-$ 控制，附参与者层面 bootstrap 95\% CI）。对每项指标，若预先指定的 Shapiro 检验通过，则采用配对 $t$ 检验（效应量为 Cohen's $d_z$）；否则采用 Wilcoxon 符号秩检验（效应量为秩二列相关 $r_{rb}$）；检验统计量在正文中报告}
\label{tab:study_typing}
\centering
\footnotesize
\setlength{\tabcolsep}{4pt}
\begin{tabular}{lccc}
\toprule
指标 & 控制 & 掩蔽 & $\Delta$ [95\% CI] \\
\midrule
准确率（\%） & -- & -- & -- [--, --] \\
WPM & -- & -- & -- [--, --] \\
CPM & -- & -- & -- [--, --] \\
尝试字符数 & -- & -- & -- [--, --] \\
首次按键延迟（ms） & -- & -- & -- [--, --] \\
\bottomrule
\end{tabular}
\end{table}

\begin{table}[!t]
\caption{六种显示条件下的主观评分（$N=$ XX；1--5 级 Likert 量表，越高越好；均值 $\pm$ 标准差）。稳定性为反向闪烁条目，舒适度为即时视觉舒适度条目（\S\ref{sec:user_study}）；Friedman 检验和 Holm 校正的成对 Wilcoxon 检验统计量在正文中报告}
\label{tab:study_ratings}
\centering
\footnotesize
\setlength{\tabcolsep}{4pt}
\begin{tabular}{lccc}
\toprule
条件 & 可读性 & 稳定性 & 舒适度 \\
\midrule
无掩蔽锚点 & -- & -- & -- \\
$n{=}2$ 掩蔽加噪声 & -- & -- & -- \\
$n{=}3$ 掩蔽加噪声 & -- & -- & -- \\
$n{=}4$ 掩蔽加噪声 & -- & -- & -- \\
$n{=}4$ 仅掩蔽 & -- & -- & -- \\
可读性优先完整配置（$n{=}4$） & -- & -- & -- \\
\bottomrule
\end{tabular}
\end{table}

\textbf{结果}：表~\ref{tab:study_typing} 和图~\ref{fig:study_typing} 汇总了打字任务。与控制条件相比，可读性优先完整配置使平均准确率变化 XX 个百分点（95\% CI [XX, XX]，$p=$ XX，效应量 XX），WPM 变化 XX（95\% CI [XX, XX]，$p=$ XX，效应量 XX），首次按键延迟变化 XX\,ms（95\% CI [XX, XX]，$p=$ XX，效应量 XX）；CPM 和尝试字符数对应的 $p$ 值分别为 XX 和 XX。表~\ref{tab:study_ratings} 和图~\ref{fig:study_ratings} 汇总了主观评分：可读性、稳定性和即时视觉舒适度的 Friedman 检验结果分别为 $\chi^2=$ XX（$p=$ XX）、$\chi^2=$ XX（$p=$ XX）和 $\chi^2=$ XX（$p=$ XX）；可读性优先完整配置与无掩蔽锚点在三个条目上的 Holm 校正成对比较 $p$ 值分别为 XX、XX 和 XX。

\begin{figure}[!t]
\centering
\includegraphics[width=\columnwidth]{figures/study_typing_paired.pdf}
\caption{控制条件与可读性优先掩蔽条件下的被试内打字表现：每条灰线代表一名参与者（每个条件下两次试验的均值）；标记表示组均值及参与者层面的 bootstrap 95\% 置信区间。各面板依次显示转录准确率、每分钟词数和首次按键延迟。}
\label{fig:study_typing}
\end{figure}

\begin{figure}[!t]
\centering
\includegraphics[width=\columnwidth]{figures/study_ratings_likert.pdf}
\caption{六种评分显示条件下的可读性、稳定性和即时视觉舒适度的 1--5 级 Likert 平均评分（误差线：标准差），条件从无掩蔽锚点排列至可读性优先完整配置。所有条目均为分值越高越好。}
\label{fig:study_ratings}
\end{figure}

}

% ----------------------------------------------------------
\subsection{补充检测与跟踪诊断}

目标检测与跟踪诊断在补充材料中报告，不纳入主要证据链。这些实验只使用一个真机拍摄位置、有限的 COCO 子集、MOT 静止帧拍摄和近似仿真指标，因此属于诊断性压力测试，不能扩展主要的传统 OCR 主张。主文省略归一化的跨任务汇总图，因为 OCR 字符恢复率、平均精度均值（mAP）、高阶跟踪准确率（HOTA）和身份识别 F1（IDF1）并非可直接比较的测量。

% ----------------------------------------------------------
\subsection{软件仿真补充实验}

本节使用软件仿真回答两个问题：（1）~经过应用渲染器后得到的单个受保护子帧，其机器可读性如何？（2）~在相机噪声为零且完全对齐时，恢复情况如何？0--2.6\% 的结果仅表征掩蔽后返回单个受保护子帧这一采集点。操作系统截图 API 可能在掩蔽前、合成后或跨已记录帧进行捕获；本研究未测试这些路径。截获原始帧的恶意软件，或对全周期进行累积的录制，都可以绕过该方法。因此，我们不提出一般性的“抗截图”主张。

\subsubsection{合成单子帧 OCR}
在 120 个合成样本上使用 Tesseract~\cite{b_smith2007}、EasyOCR~\cite{b_easyocr2020} 和 Surya~\cite{b_surya2024} 后，三个引擎的字符恢复率从 94.0--97.0\% 降至 0--2.6\%（表~\ref{tab:ocr_corpus}）。该结果仅覆盖当前语料库、参数和引擎集合；分层细分可以检查类别趋势，但不足以证明对所有语言、布局或识别器均有效。

\begin{table}[!t]
\caption{合成单子帧 OCR：三引擎字符恢复率（120 个样本，$n=4$，$\varepsilon=8/255$）}
\label{tab:ocr_corpus}
\centering
\footnotesize
\begin{tabular}{lccc}
\toprule
OCR 引擎 & 原始帧 & 单子帧 & 降幅 \\
\midrule
Tesseract & 94.0\% & 0.0\% & 94.0\% \\
EasyOCR & 94.1\% & 0.0\% & 94.1\% \\
Surya & 97.0\% & 2.6\% & 94.4\% \\
\bottomrule
\end{tabular}
\end{table}

\begin{figure}[!t]
\centering
\includegraphics[width=\columnwidth]{figures/multiengine_ocr.pdf}
\caption{三个 OCR 引擎在原始帧和单个受保护子帧上的字符恢复率；误差线为以语料样本为单位重采样所得的 95\% bootstrap 区间。}
\label{fig:multiengine_ocr}
\end{figure}

图~\ref{fig:multiengine_ocr} 将这些结果可视化。所有引擎的单子帧恢复率均不超过 2.6\%，说明该发现并非由某一个识别器驱动。更广泛的推广仍需增加引擎和语料库。

\subsubsection{理论安全边界}
\label{sec:security_frontier}
当攻击者累积覆盖全周期且完成配准的帧时，在式~(\ref{eq:noise_complement}) 的理想线性模型中，互补子帧可以重新求和。归档中逐样本最强的数字攻击恢复率为 95.2\%；纯全周期时域平均的恢复率为 94.3\%。这是所评估线性分解的重建边界，而非正式的信息论定理。在 150 帧均值下，高抑制档的汇总恢复率仍为 47.9\%；排除欠曝光位置后为 53.9\%。字符恢复率不必与点亮像素比例成比例，因为识别结果还取决于内容和物理成像路径。

\begin{table}[!t]
\caption{归档的数字重建边界（Tesseract）}
\label{tab:digital_boundary}
\centering
\footnotesize
\setlength{\tabcolsep}{3pt}
\begin{tabular}{lrrrrr}
\toprule
设置 & $N$ & 字符恢复率 & 完全匹配率 & SSIM & $\Delta E_{00}$ \\
\midrule
基础：循环均值 & 120 & 94.3\% & 85.8\% & -- & -- \\
基础：逐样本最强 & 120 & 95.2\% & 91.7\% & -- & -- \\
可读性优先：循环均值 & 16 & 87.9\% & 50.0\% & 0.948 & 1.46 \\
高抑制：循环均值 & 16 & 56.6\% & 6.3\% & 0.764 & 4.87 \\
高抑制：反色时隙 & 16 & 93.0\% & 81.2\% & -- & -- \\
\bottomrule
\end{tabular}
\vspace{1mm}

\raggedright\scriptsize 120 样本的基础攻击与 16 样本的档位消融是两项独立归档的数字评估。“逐样本最强”针对每个样本，从七种预先指定的重建中选择最佳结果。档位的 SSIM 和 $\Delta E_{00}$ 描述全周期的重建质量，而非物理质量或用户感知质量。
\end{table}

主要证据链之外还保留了其他软件代理基线、参数扫描、检测和跟踪诊断。由于其物理条件、用户体验代价和指标与真机拍摄 OCR 的估计目标并不匹配，因此不将其作为主文证据。

\deferreduserstudy

%% ============================================================
%% VI. DISCUSSION
%% ============================================================
\section{讨论}
\label{sec:discussion}

\subsection{证据层级与主张范围}

最直接获得支持的结论是固定条件下的观测。在一条 S600 链路的共同标称设置下，所评估档位降低了三个开源 OCR 引擎在原始输入上的平均恢复率。固定的三引擎预处理网格能够恢复出更多文本：可读性优先档的网格最优选择达到 40.2\%，高抑制档达到 13.7\%，因此不能将原始输入视为攻击上界。相对于未保护输入，剩余的配对差值量化的是该网格内的降幅，而非对优化攻击者的抵抗能力。该结果既不构成曝光区间，也不构成一般性的防拍摄主张。语料库不平衡、仅使用一个网络摄像头、档位顺序固定，以及缺少亮度匹配的对照，使我们无法作出跨设备推断或机制隔离推断。VLM、积分和仿真结果用于界定边界，而非扩展主要主张。被动基线属于软件代理，不能证明本方法优于物理滤镜或既有时序显示系统。

\subsection{线性积分的数学张力}

该方案利用了人类时序感知与相机短曝光之间的差异。然而，其数学边界建立在理想的线性屏幕--相机辐亮度模型上，而不是将人类视觉视为线性求和算子。在该模型中，互补子帧可以重新求和以恢复原图（$\sum\mathbf{I}_k = \mathbf{I}$），因此能够累积一个完整配准循环的相机会击破这一线性分解。关于 VLM 对抗鲁棒性的研究~\cite{b_jiang2025} 进一步表明，大模型能够从受扰图像中提取部分语义信息；\S\ref{sec:vlm} 的探针实验验证了这一点。VLM 能够从单张短曝光照片中恢复大字号内容，并从 150 帧视频均值和长曝光视图中获得更强信号。在所有已评估重建中，归档数字评估的逐样本最强攻击达到 95.2\%；纯全周期时域均值的恢复率达到 94.3\%（\S\ref{sec:security_frontier}）。

\subsection{物理层混杂因素与仿真--实拍差距}

仿真--实拍差距与若干物理层效应相容，但当前设计未能单独识别其中任何一种效应。

\textbf{方向 1：实拍短曝光恢复率高于仿真（有利于攻击者）。} 仿真单子帧 OCR 的字符恢复率为 0--2.6\%（表~\ref{tab:ocr_corpus}），而真机拍摄在相近掩蔽条件下的短曝光恢复率为 19.2\%（仅掩蔽档，表~\ref{tab:real_ocr_full_pool}）。两种物理效应可能对此有所贡献。第一，滚动快门行混合可能使不同行积分不同的子帧内容，但本研究未测量 S600 的读出时间。第二，LCD 响应转换可能混合相邻子帧，但未表征面板的灰阶转换行为。真机拍摄测量包含这两种效应，但本实验无法识别它们各自的贡献。因此，仿真是理想化的单子帧参照，而非物理下界。

\textbf{方向 2：实拍积分恢复率低于仿真。} 逐样本最强数字攻击的恢复率为 95.2\%（\S\ref{sec:security_frontier}），而实拍 150 帧均值在可读性优先档和高抑制档上的恢复率分别为 71.1\% 和 47.9\%。共同设置下的数值分别为 79.7\% 和 53.9\%。ISP 处理、moir\'{e}（摩尔纹）、模糊和不完美配准都可能有所贡献，但归档无法识别它们是否存在及其幅度。专门设计的积分攻击因而可能缩小观测到的差距。

因此，仿真和物理实验回答的是不同问题。仿真评估理想化的单子帧变换；物理恢复率则测量包含未解决混杂因素在内的完整归档屏幕--相机流水线。

\subsection{探索性高抑制档的作用与代价}

探索性高抑制档通过条纹和字形扰动破坏严格完备性，并与较低的积分恢复率相关。其全周期数字重建 SSIM 为 0.764，可读性优先档则为 0.948；$\Delta E_{00}$ 分别为 4.87 和 1.46（表~\ref{tab:digital_boundary}）。同一数字消融暴露出直接的单槽失效：经过对比度拉伸后，反色帧的平均字符恢复率为 93.0\%，完全匹配率为 81.2\%。这一观测到的 93.0\% 最佳数字恢复率是该档位最强的已观测数字攻击，因此 56.6\% 的时域平均结果并非最坏情形安全界限。在共同设置池中，高抑制档的短曝光字符恢复率为 5.6\%，150 帧均值下为 53.9\%。它没有达到描述性的 $<5\%$ 目标，且只能缓解积分泄漏。用户接受度尚未测试。

这一档位区分是解释结果的核心。可读性优先档是唯一与拟开展舒适度研究相联系的档位，但它在显式字段和积分攻击下更弱。高抑制档在若干视图中降低了观测恢复率，但可利用的反色时隙和缺失的可用性证据使其不能支持一般性的安全主张。因此，本文支持的是一个有限的可行性空间，而非部署建议。

\subsection{VLM 攻击者与对策方向}
\label{sec:vlm_defense}

VLM 的高恢复率集中在大字号短片段上（\S\ref{sec:vlm}），而时域平均能够提供更强信号。这一模式与保留的字形轮廓和语言先验补全相符，但本实验未单独识别其中任何一种机制。我们保留该结果，是因为它限制了方法的适用范围：降低传统 OCR 恢复率并不意味着能够抵抗 VLM。由此可以提出三个尚未验证的假设：字体大小自适应粒度、字形诱饵，以及字段级渲染策略。互补诱饵在理想的全周期积分下仍会抵消，且这些设想均未经过物理验证。因此，实际目标应表述为提高攻击成本，而不是使可见信息无法恢复。

\subsection{持续视频攻击者}

所评估的视频攻击并未证明四帧即可绕过。拍摄流水线以标称 60\,fps 记录 150 个短曝光帧，报告的时域均值图像对完整连拍取平均，对应约 2.5\,s 的标称拍摄时长。由于持续平均会累积多个显示循环，该攻击不需要掩蔽种子、相位或精确的帧选择。其汇总恢复率在可读性优先档上为 71.1\%，在高抑制档上为 47.9\%；排除欠曝光位置后，分别升至 79.7\% 和 53.9\%。归档还包含从同一完整连拍中选取的聚合视图，但它们无法证明从独立采集的短窗口中也能恢复。时长--恢复率曲线和预先指定的短窗口分析尚未评估。因此，应区分单次拍摄适用范围与已测得的持续录制边界，同时不指定缺乏证据支持的最短绕过时长。

\subsection{部署路径}

真实部署需要驱动级帧缓冲注入、面板/背光协调和经过验证的时序。这些步骤会引入色彩管理、亮度饱和、丢帧和同步误差。在匹配短曝光分析中，可读性优先档的显式字段恢复率仍为 17.7\%，且在积分攻击下字符恢复率较高。因此，凭据保护不应仅依赖该档位。候选的字段级缓解措施需要分别进行可读性、曝光、视频和 VLM 评估。

\subsection{所需的验证性物理协议}

验证性采集应在每个内容--几何条件区块内随机化档位顺序，并记录墙钟时间批次标识符。实验应锁定并记录曝光模式、增益、白平衡、焦点、编解码器、固定的显示器亮度设置和环境照度。应使用光电二极管或高速传感器测量时隙时序和显示器物理亮度。每种时序档位都应配备在时间平均亮度和叠加外观上匹配的静态对照。这些对照可以将档位关联与时间漂移、占空比亮度、静态叠加、面板响应和相机处理区分开来。在获得此类数据之前，现有归档无法支持有关时序机制的因果主张。

\subsection{伦理声明}

本研究具有防御性质，旨在保护屏幕内容免遭未经授权的拍摄。攻击分析（\S\ref{sec:security_frontier}）旨在诚实评估安全边界，而非提供攻击指南。拟开展的用户研究属于最小风险的人体参与研究；目前尚未收集任何参与者数据。只有在具备权限的机构或院系主管部门出具书面批准、豁免或免审认定后，才会开始招募。协议规定了知情同意、光敏感筛查和自愿退出，详见 \S\ref{sec:user_study}。

\subsection{局限性}

\begin{enumerate}
\item 用户研究（\S\ref{sec:user_study}）目前是一项协议，而非已完成的证据：尚无参与者数据，因此可读性优先档的可用性仍未得到验证。即使在受控的 240\,Hz 实验室中完成研究，外部有效性仍需通过更广泛的参与者样本和显示条件加以检验。
\item 即时视觉舒适度是在短时暴露后采用单条目量表测得，不能推广至长期疲劳或临床视觉不适。
\item 该概念验证（PoC）仅位于应用层，不包含操作系统驱动程序或硬件背光。亮度、高动态范围（HDR）和色差结论来自数字模型，缺少光度学或面板响应测量。实际播放时隙时序未通过光电二极管、高速相机或示波器验证。由于实验缺少亮度匹配的静态对照和匹配的无噪声增强档位，因此无法将时序稀疏性、占空比亮度、静态叠加、面板响应、相机曝光和 ISP 处理作为机制区分开来。
\item 真机拍摄 OCR 仅使用一个 S600 UVC 网络摄像头。显示器亮度设置和采集相位保持固定，但未记录显示器物理亮度和环境照度。传感器增益、自动曝光语义、白平衡状态和焦点状态均未保留。档位顺序固定而非随机。因此，时间、批次和相机自动行为可能混杂档位或几何条件比较，但无法重建其方向和幅度。重复内容和拍摄重采样汇总不能替代跨设备复现。
\item 真机拍摄检测在单一位置使用了 150 张 COCO 图像；跟踪使用来自一个 MOT17 序列的 450 次静止帧拍摄，而非连续视频。仿真检测使用 8~张图像和近似指标。
\item VLM 评估覆盖一个相机、0$^{\circ}$ 正轴、一个档位和三个商业 API 模型（每个单元格 12 个视频片段）。模型可能随时间更新。公开文档内容可能通过训练数据记忆抬高恢复率。我们未评估多循环自适应重建或掩蔽前恶意软件截获。全周期积分仍是所评估线性设计的重建边界。
\item 传统 OCR 评估覆盖 Surya、EasyOCR 和 Tesseract，但不包含商业 OCR API。预处理压力测试穷举了针对这些引擎预先声明的六输入网格，而非连续参数调优、学习式增强、图像复原或识别器特定优化。在所评估集合内，最优引擎与输入的联合选择有利于攻击者，但 Google Cloud Vision、Azure AI Vision、Baidu OCR 或未来识别器仍可能得到更高恢复率。
\item 基线比较使用未匹配的软件代理，而非经过物理测试的隐私滤镜或既有时序显示系统。这些比较用于提供观测权衡的背景，但不能证明本方法相对于这些替代方案具有比较优势。
\item 可读性优先档不适合作为敏感字段的唯一保护层：匹配短曝光条件下的字段微平均完整恢复率为 17.7\%，积分/VLM 攻击能恢复更多内容。仍需独立的访问控制。
\item 由 12 个内容项构成的语料库并非全英文/ASCII，而是混合语料：CET6 页面和两个混合语言片段包含中文字符。该语料库规模过小且经过选择性构造，无法支持文字系统层面的比较；专门的非拉丁文字评估仍属未来工作。
\end{enumerate}


%% ============================================================
%% VII. CONCLUSION
%% ============================================================
\section{结论}
\label{sec:conclusion}

我们在一个标称 UVC 控制设置下对显示处理进行了档位级测量。在 8 种共同设置几何条件下，每个档位使用相同的 288 个匹配单元（12 个内容项 $\times$ 8 种几何条件 $\times$ 3 次重复），未保护条件、可读性优先档和高抑制档的平均恢复率分别为 94.5\%、17.8\% 和 5.6\%。按内容项聚类的配对差异分别为 76.7 个百分点（[74.3, 79.2]）和 88.9 个百分点（[85.5, 91.9]）。这些结果支持三个开源 OCR 引擎在固定链路上的关联。它们既未确定曝光区间，也未确立时序机制的因果关系。可读性优先档仍有 17.7\% 的匹配字段微平均完整恢复率和 95.0\% 的 P99 字符恢复率，而高抑制档未达到其描述性的 $<5\%$ 目标。

绕过证据同样构成严格限制。150 帧时域均值在可读性优先档和高抑制档上的字符恢复率分别为 71.1\% 和 47.9\%；共同设置下的数值分别为 79.7\% 和 53.9\%。在完整的 9 几何条件敏感性池中，标称控制设置下的长曝光恢复率在可读性优先档上为 60.9\%，未保护条件为 47.3\%。正轴 VLM 探针也恢复了大量大字号内容。在完整固定网格上，三引擎网格最优选择在可读性优先档和高抑制档上分别达到 40.2\% 和 13.7\%。不受约束或学习式预处理以及商业 OCR 均不在当前证据所支持的条件范围内。

因此，本研究贡献的是一项有边界的物理测量及明确的失效约束，而非可部署的防相机方案。下一步最关键的实验是采用匹配静态对照的随机化采集、曝光--恢复率曲线、可用性测试和跨设备复现。

%% ============================================================
%% DATA AND CODE AVAILABILITY
%% ============================================================
\section*{数据与代码可用性}

软件、实验配置、去标识化的逐样本 OCR 结果、VLM 记录和分析脚本维护于 \url{https://github.com/docandyc/privacy-display}。该仓库目前仍可变更，不作为不可变的投稿产物提供。为支持复现，定稿版本将提供冻结的发布标签或归档 DOI。VLM 记录包含端点标识符、时间戳、提示词、解码参数、输入哈希或路径、解析状态和原始输出。归档标识符 \texttt{capture\_hardened} 和 \texttt{vlm} 均表示高抑制档。若用户研究得以开展，可识别身份的注册数据将不予公开；仅公开去标识化的结果和审计信息。

%% ============================================================
%% ACKNOWLEDGMENT
%% ============================================================
\section*{致谢}
\TODO{致谢待填写。}

OpenAI Codex~\cite{b_openai_codex2026} 协助开发分析脚本、执行一致性检查，并对摘要、引言、实验评估、讨论和结论进行英文修订。作者确定科学决策，依据归档数据核验生成的分析，审阅所有保留的文本与代码，并对本文承担责任。


%% ============================================================
%% REFERENCES
%% ============================================================
\bibliographystyle{IEEEtran}
\bibliography{refs}
\EOD

\vfill

\end{document}
